% !TeX program = pdflatex
% !BIB program = biber
\documentclass[a4paper,10pt,twoside,openany]{book}
% Package order is intentional; load cross-reference packages last.
\usepackage[T1]{fontenc}
\usepackage[utf8]{inputenc}
\usepackage[english]{babel}
\usepackage{amsmath,amssymb,amsthm,mathtools,bm}
\usepackage{mathpazo}
\usepackage{microtype}
\usepackage[a4paper,inner=26mm,outer=22mm,top=24mm,bottom=25mm,
  headheight=14pt,headsep=7mm,footskip=12mm]{geometry}
\usepackage{xcolor}
\usepackage{graphicx}
\usepackage{booktabs,longtable,array,tabularx}
\usepackage{siunitx}
\usepackage{enumitem}
\usepackage{fancyhdr}
\usepackage{titlesec}
\usepackage{caption}
\usepackage{tikz,pgfplots}
\usetikzlibrary{arrows.meta,calc,positioning,patterns}
\pgfplotsset{compat=1.18}
\usepackage{makeidx}
\makeindex
\usepackage{csquotes}
\usepackage[backend=biber,style=numeric-comp,sorting=none,
  giveninits=true,maxbibnames=99,doi=true,isbn=true,url=false]{biblatex}
\addbibresource{bibliography/references.bib}
\usepackage{hyperref}
\usepackage{bookmark}
\usepackage[nameinlink,noabbrev]{cleveref}

% Semantic notation. Do not redefine TeX accents or standard commands.
\newcommand{\ii}{\mathrm{i}}
\newcommand{\ee}{\mathrm{e}}
\newcommand{\dd}{\mathop{}\!\mathrm{d}}
\newcommand{\vect}[1]{\bm{#1}}
\newcommand{\mat}[1]{\mathbf{#1}}
\newcommand{\op}[1]{\hat{#1}}
\newcommand{\adj}{^{\dagger}}
\DeclareMathOperator{\Tr}{Tr}
\DeclareMathOperator{\diag}{diag}
\DeclareMathOperator{\rank}{rank}
\DeclareMathOperator{\Cov}{Cov}
\DeclareMathOperator{\Var}{Var}
\DeclareMathOperator{\E}{\mathbb{E}}
\DeclareMathOperator*{\argmin}{arg\,min}
\DeclarePairedDelimiter{\norm}{\lVert}{\rVert}
\DeclarePairedDelimiter{\abs}{\lvert}{\rvert}
\newcommand{\ket}[1]{\lvert #1\rangle}
\newcommand{\bra}[1]{\langle #1\rvert}
\newcommand{\braket}[2]{\langle #1\vert #2\rangle}
\newcommand{\gammabar}{\bar{\gamma}}
\newcommand{\Mplus}{M_{+}}
\newcommand{\Mminus}{M_{-}}
\newcommand{\Btx}{B_{1}^{+}}
\newcommand{\TtwoStar}{T_{2}^{*}}
\newcommand{\RtwoStar}{R_{2}^{*}}
\newcommand{\TR}{\mathrm{TR}}
\newcommand{\TE}{\mathrm{TE}}
\newcommand{\TI}{\mathrm{TI}}
\newcommand{\FOV}{\mathrm{FOV}}
\newcommand{\SNR}{\mathrm{SNR}}
\newcommand{\CNR}{\mathrm{CNR}}
\newcommand{\RF}{\mathrm{RF}}
\newcommand{\FT}{\mathcal{F}}
\newcommand{\sinc}{\operatorname{sinc}}
% sinc(x) = sin(pi x)/(pi x), including its continuous value at zero.
% Always use gamma for rad/(s T), gammabar for Hz/T; never interchange.

\definecolor{MRIaccent}{HTML}{254B63}
\definecolor{MRIgray}{HTML}{505B62}
\hypersetup{
  pdftitle={MRI Physics: From Quantum Spin Dynamics to Advanced Acquisition and Reconstruction},
  pdfauthor={},pdfsubject={Graduate MRI physics reference: foundations edition},
  colorlinks=true,linkcolor=MRIaccent,citecolor=MRIaccent,urlcolor=MRIaccent,
  bookmarksnumbered=true,linktoc=all
}
\setcounter{secnumdepth}{2}
\setcounter{tocdepth}{2}
\setlength{\parindent}{1.1em}
\setlength{\parskip}{0pt plus .4pt}
\setlength{\emergencystretch}{1em}
\raggedbottom
\setlist{nosep}
\sisetup{per-mode=symbol,detect-all}
\captionsetup{font=small,labelfont=bf,margin=8mm}
\titleformat{\chapter}[display]
  {\normalfont\huge\bfseries\color{MRIaccent}}
  {\normalfont\large\scshape\chaptertitlename\ \thechapter}{8pt}{}
\titlespacing*{\chapter}{0pt}{8pt}{24pt}
\titleformat{\section}{\normalfont\large\bfseries}{\thesection}{.7em}{}
\titleformat{\subsection}{\normalfont\normalsize\bfseries}{\thesubsection}{.7em}{}
\pagestyle{fancy}
\fancyhf{}
\fancyhead[LE,RO]{\small\thepage}
\fancyhead[RE]{\small\nouppercase{\leftmark}}
\fancyhead[LO]{\small\nouppercase{\rightmark}}
\renewcommand{\headrulewidth}{.3pt}
\renewcommand{\chaptermark}[1]{\markboth{\thechapter\quad #1}{}}
\renewcommand{\sectionmark}[1]{\markright{\thesection\quad #1}}
\fancypagestyle{plain}{\fancyhf{}\fancyfoot[C]{\small\thepage}
  \renewcommand{\headrulewidth}{0pt}}
% Unboxed pedagogical environments; examples and results share one counter.
\newtheoremstyle{mriprosemath}{8pt}{6pt}{\normalfont}{}{\bfseries}{.}{.5em}{}
\theoremstyle{mriprosemath}
\newtheorem{definition}{Definition}[chapter]
\newtheorem{physicalinterpretation}[definition]{Physical Interpretation}
\newtheorem{derivation}[definition]{Derivation}
\newtheorem{workedexample}[definition]{Worked Example}
\newtheorem{importantresult}[definition]{Important Result}
\newtheorem{commonmisconception}[definition]{Common Misconception}
\newtheorem{clinicalconnection}[definition]{Clinical Connection}
\newtheorem{researchnote}[definition]{Research Note}
\newtheorem{challengeproblem}{Challenge Problem}[chapter]
\crefname{definition}{definition}{definitions}
\crefname{challengeproblem}{problem}{problems}
\Crefname{challengeproblem}{Problem}{Problems}
\newenvironment{solution}[1]
  {\par\medskip\noindent\textbf{Solution to \Cref{#1}.}\ \ignorespaces}
  {\par\medskip}
% Figure vocabulary: RF/acquisition blue; gradients gray; annotations black.
\tikzset{
  mri axis/.style={-{Stealth[length=2mm]},thin},
  mri rf/.style={MRIaccent,thick},
  mri gradient/.style={MRIgray,thick},
  mri signal/.style={MRIaccent,thick},
  mri guide/.style={gray,densely dashed,thin},
  mri label/.style={font=\small}
}
\pgfplotsset{every axis/.append style={
  axis line style={thin},tick label style={font=\small},
  label style={font=\small},legend style={font=\small,draw=none},
  grid style={gray!20},width=.8\textwidth,height=.42\textwidth}}

% For focused drafting, uncomment one or more files (no extension):
% \includeonly{chapters/03_classical_magnetization_bloch}
\begin{document}
\frontmatter
\begin{titlepage}
\thispagestyle{empty}
\vspace*{27mm}
{\noindent\color{MRIaccent}\rule{\textwidth}{1.1pt}\par}
\vspace{12mm}
{\fontsize{38}{43}\selectfont\bfseries MRI Physics\par}
\vspace{9mm}
{\Large From Quantum Spin Dynamics to\par
Advanced Acquisition and Reconstruction\par}
\vfill
{\small\scshape Graduate reference textbook\par}
\vspace{3mm}
{\small Foundations edition\par}
\vspace{10mm}
{\noindent\color{MRIaccent}\rule{\textwidth}{.4pt}\par}
\end{titlepage}
\chapter{Scope and Reading Structure}
\label{front:scope}
This project is an advanced reference for a reader with graduate training
in applied and medical physics. Its organizing argument runs from spin
dynamics to measurable signals, acquisition design and inference.
Prerequisites are developed where they support that argument.

This edition completes the foundational block in Chapters~1--6:
mathematical tools, quantum spin physics, classical magnetization,
RF excitation, spatial encoding and Fourier signal formation.
These chapters contain derivations, numerical examples, vector figures
and thirty challenge problems with worked solutions in Appendix~\ref{app:e}.
Chapters~7--27 and the remaining reference appendices retain their planned
hierarchy and are not yet completed exposition. No chapter is assigned
an artificial page budget.

The default reading order is the printed chapter order. The EPG treatment
in \cref{ch:13} supplies a second, computational reading of the echo trains
introduced in \cref{ch:08} and the transient models in \cref{ch:12}.
General estimation ideas start in \cref{ch:01}, become measurement models
in \cref{ch:11}, and become reconstruction algorithms in \cref{ch:16}.
This avoids requiring the entire inverse-problem chapter before the
first quantitative experiment.

\section*{Reference framework}
The initial reference database anchors sequence design in
Bernstein, King and Zhou \cite{bernstein2004} and spin dynamics in
Levitt \cite{levitt2008}. Method-specific anchors are the
Stejskal--Tanner diffusion experiment \cite{stejskal1965},
SENSE \cite{pruessmann1999}, GRAPPA \cite{griswold2002},
compressed-sensing MRI \cite{lustig2007}, magnetic resonance
fingerprinting \cite{ma2013}, and the EPG review by
Weigel \cite{weigel2015}. The foundational chapters additionally cite primary work on induction,
spin echoes, diffusion, exchange, encoding, sampling and reception.
References attached to later chapter outlines identify intended sources,
not completed treatments.

\tableofcontents
% This is front-matter content, deliberately input after \begin{document}.
\chapter{Notation and Mathematical Conventions}
\label{front:conventions}
The definitions below form a contract for the manuscript, figures and
numerical examples. They fix conventions rather than substitute for the
derivations in \cref{ch:01,ch:02,ch:03,ch:04,ch:05,ch:06,ch:13}.
SI units are used unless another unit is explicitly shown. Angles in
exponentials are in radians. Spatial coordinates are right-handed scanner
coordinates; patient orientation must be specified separately.

\section*{Spin, precession and RF phase}
For proton imaging, the gyromagnetic ratio is positive:
$\gamma>0$ in $\si{\radian\per\second\per\tesla}$ and
$\gammabar=\gamma/(2\pi)$ in $\si{\hertz\per\tesla}$.
The static field is $\vect B_0=B_0\vect e_z$, with $B_0>0$, and
$\omega_0=\gamma B_0$, $f_0=\gammabar B_0$ denote positive frequency
magnitudes. The magnetic moment and Zeeman Hamiltonian are
\begin{equation}
 \op{\vect\mu}=\gamma\op{\vect S},\qquad
 \op H=-\op{\vect\mu}\cdot\vect B,\qquad
 \op S_j=\frac{\hbar}{2}\sigma_j .
 \label{eq:convention:hamiltonian}
\end{equation}
Here $\sigma_j$ are Pauli matrices and $\hbar$ is the reduced Planck constant.
We use $\dot{\vect M}=\gamma\vect M\times\vect B$ for torque-only evolution.
Consequently $\Mplus=M_x+\ii M_y$ evolves as
\begin{equation}
 \Mplus(t)=\Mplus(0)\ee^{-\ii\omega_0t}
 \quad\text{in a uniform static field without relaxation}.
 \label{eq:convention:precession}
\end{equation}
Positive $\omega_0$ therefore does not mean a positive right-handed
rotation angle about $+z$.

The rotating-frame transverse variable is
$\Mplus^{\mathrm{rot}}(t)=\ee^{+\ii\omega_{\mathrm{RF}}t}
\Mplus^{\mathrm{lab}}(t)$, where $\omega_{\mathrm{RF}}>0$ is the RF carrier
magnitude. Detuning is $\Delta\omega=\gamma B_z-\omega_{\mathrm{RF}}$;
free rotating-frame evolution has factor $\ee^{-\ii\Delta\omega t}$.
For a physical RF field written
$\vect B_{\mathrm{RF}}(t)=\operatorname{Re}\{
  \vect b(t)\ee^{-\ii\omega_{\mathrm{RF}}t}\}$,
define its co-rotating component by
\begin{equation}
 \Btx(t)=\frac{b_x(t)+\ii b_y(t)}{2}.
 \label{eq:convention:b1plus}
\end{equation}
The phasors $b_x,b_y$ and $\Btx$ have units of tesla. With this
negative-time-exponential convention, a clockwise circular field has
$b_y=-\ii b_x$ and $\Btx=b_x$. A linearly polarized peak field
$B_{\mathrm{pk}}\cos(\omega_{\mathrm{RF}}t)\vect e_x$ has
$\Btx=B_{\mathrm{pk}}/2$. The symbol $B_1$ alone must be qualified as
a physical amplitude, phasor amplitude or co-rotating amplitude.

Let $R_{\vect n}(\theta)$ denote an active right-handed rotation.
An on-resonance hard pulse with constant phase
$\phi=\arg\Btx$, negligible relaxation and flip magnitude
$\alpha=\gamma\int\abs{\Btx(t)}\dd t$ acts as
$R_{\vect n}(-\alpha)$, with
$\vect n=(\cos\phi,\sin\phi,0)$.
Thus a positive $90^\circ$ pulse about $+x$ sends $+z$ to $+y$.
The corresponding spin propagator is
$U=\exp(+\ii\alpha\,\vect n\cdot\vect\sigma/2)$.
This integral flip-angle rule is not asserted for arbitrary phase
modulation or detuning.

\section*{Fourier transforms, reception and encoding}
The continuous spatial transform and its inverse are
\begin{align}
 \widetilde f(\vect k)
 &=\int_{\mathbb R^d}f(\vect r)
       \ee^{-\ii2\pi\vect k\cdot\vect r}\dd^d\vect r,
 \label{eq:convention:fourier-forward}\\
 f(\vect r)
 &=\int_{\mathbb R^d}\widetilde f(\vect k)
       \ee^{+\ii2\pi\vect k\cdot\vect r}\dd^d\vect k.
 \label{eq:convention:fourier-inverse}
\end{align}
Here $\vect r$ is in metres and $\vect k$ is in cycles per metre.
The same negative sign defines the temporal transform
$\widetilde s(f)=\int s(t)\ee^{-\ii2\pi ft}\dd t$.
A freely precessing component $\ee^{-\ii2\pi\Delta f t}$ therefore
appears at transform coordinate $f=-\Delta f$; a displayed spectroscopic
axis increasing with physical resonance offset must explicitly reverse
that coordinate. The normalized sinc is
$\sinc u=\sin(\pi u)/(\pi u)$.

For an interval without RF refocusing and with
$B_z(\vect r,t)=B_0+\vect G(t)\cdot\vect r$,
\begin{equation}
 \vect k(t)=\vect k(t_0)+\gammabar\int_{t_0}^{t}\vect G(\tau)\dd\tau .
 \label{eq:convention:kspace}
\end{equation}
$\vect G$ is in $\si{\tesla\per\metre}$, while
$\dot{\vect G}$ is in $\si{\tesla\per\metre\per\second}$.
Across RF refocusing, transverse conjugation and RF phase must be
propagated explicitly; a toggling-frame gradient integral is valid only
for the stated coherence pathway.

Receiver phase is chosen so that the demodulated signal is linear in
$\Mplus^{\mathrm{rot}}$. Define $c_c(\vect r)$ as the complex receive
weight multiplying this variable for coil $c$; any conjugation in a
reciprocity-field convention is absorbed in this definition.
Under stationary, linear-gradient, time-independent-object and
negligible-intrareadout-decay assumptions, the signal reduces to
\begin{equation}
 s_c(t)=\int c_c(\vect r)\,m(\vect r)
       \ee^{-\ii2\pi\vect k(t)\cdot\vect r}\dd^3\vect r .
 \label{eq:convention:ideal-signal}
\end{equation}
$m$ is the prepared complex transverse magnetization density, with
fixed receiver scaling absorbed into $c_c$. It is not automatically
proton density. Off-resonance, relaxation, motion and chemical species
require the more general model in \cref{ch:06}.

The one-dimensional DFT uses
$X_p=\sum_{n=0}^{N-1}x_n\ee^{-\ii2\pi pn/N}$ and
$x_n=N^{-1}\sum_{p=0}^{N-1}X_p\ee^{+\ii2\pi pn/N}$.
Multidimensional inverse normalization is the product of dimension
lengths. A unitary DFT must be marked explicitly. Shifting the array
origin changes indexing, not transform signs. For uniform Cartesian
sampling, $\FOV=1/\Delta k$ and the nominal grid spacing is
$\Delta x=\FOV/N$; effective resolution depends on the acquired support,
weighting, motion and reconstruction.

\section*{EPG state and noise conventions}
For a dimensionless spatial phase coordinate $\psi\in[0,2\pi)$,
define signed-order EPG coefficients by
\begin{equation}
 F_n^\pm=\frac{1}{2\pi}\int_0^{2\pi}M_\pm(\psi)
             \ee^{-\ii n\psi}\dd\psi,\qquad
 Z_n=\frac{1}{2\pi}\int_0^{2\pi}M_z(\psi)
             \ee^{-\ii n\psi}\dd\psi ,
 \label{eq:convention:epg}
\end{equation}
where $M_-=M_x-\ii M_y$, $n\in\mathbb Z$,
$F_n^-=(F_{-n}^+)^*$ and $Z_{-n}=Z_n^*$.
The observable uniform-voxel transverse average is $F_0^+$.
A gradient phase $\Mplus\mapsto\Mplus\ee^{-\ii q\psi}$ with
integer $q$ implies
$F_n^+\mapsto F_{n+q}^+$ and $F_n^-\mapsto F_{n-q}^-$,
where the right-hand sides are pre-gradient states.
An implementation using only nonnegative orders must state how it
folds conjugate states. RF matrices in \cref{ch:13} must follow the
rotation sign above, not be copied without translation.

Proper complex Gaussian noise is denoted
$\vect n\sim\mathcal{CN}(0,\mat\Psi)$ with
$\mat\Psi=\E[\vect n\vect n\adj]$ and
$\E[\vect n\vect n^{\mathsf T}]=0$.
For one channel with independent real and imaginary variance $\sigma^2$,
$\E[\abs n^2]=2\sigma^2$. This distinction must accompany SNR and
likelihood formulas.

\section*{Notation register}
\renewcommand{\arraystretch}{1.15}
\begin{longtable}{@{}p{.19\textwidth}p{.51\textwidth}p{.23\textwidth}@{}}
\toprule Symbol & Meaning & Unit or convention\\ \midrule
\endfirsthead
\toprule Symbol & Meaning & Unit or convention\\ \midrule
\endhead
$\vect r,\vect k$ & Position; spatial frequency & m; cycles/m\\
$t,\omega,f$ & Time; angular frequency; frequency & s; rad/s; Hz\\
$\gamma,\gammabar$ & Angular and cyclic gyromagnetic ratios & rad/(s\,T); Hz/T\\
$\vect B,\Btx$ & Magnetic flux density; co-rotating RF phasor & T\\
$\vect G,\dot{\vect G}$ & Gradient strength; slew rate & T/m; T/(m\,s)\\
$\vect M,M_0$ & Magnetization; equilibrium magnitude & A/m\\
$\vect p,p_0$ & Bloch polarization vector; equilibrium polarization & Dimensionless\\
$T_{\mathrm{bath}},k_{\mathrm B}$ & Bath temperature; Boltzmann constant & K; J/K\\
$\omega_1,\Omega$ & RF nutation rate; effective rotation rate & rad/s\\
$\Delta\omega,\sigma_\omega$ & Detuning; angular-offset standard deviation & rad/s\\
$T_p,\Delta t$ & RF pulse duration; sample dwell time & s\\
$L,K,\Delta V$ & FOV length; aperture half-width; voxel volume & m; cycles/m; $\mathrm m^d$\\
$\rho_s$ & Spin number density & $\si{\per\cubic\metre}$\\
$\op\rho$ & Density operator, distinct from spin density & Trace one\\
$T_1,T_2,\TtwoStar$ & Longitudinal, transverse, effective decay times & s\\
$R_1,R_2,\RtwoStar$ & Corresponding inverse times & $\si{\per\second}$\\
$\alpha,\phi$ & RF flip magnitude and phase & rad\\
$\TR,\TE,\TI$ & Repetition, echo and inversion times & s\\
$c_c,s_c$ & Receive weight and complex signal, coil $c$ & Calibration-dependent\\
$\mat A,\mat A\adj$ & Encoding matrix and its Hermitian adjoint & Model-dependent\\
$\mat\Psi,\sigma^2$ & Complex covariance; real-channel variance & Signal units squared\\
$D,\mat D$ & Diffusivity; diffusion tensor & $\si{\square\metre\per\second}$\\
$b,\mat B_b$ & Diffusion weighting and b-matrix & $\si{\second\per\square\metre}$\\
$\vect q$ & Angular diffusion wavevector & rad/m; not $\vect k$\\
$\mathcal M_j$ & Gradient moment of order $j$ & $\mathrm{T\,s^{j+1}/m}$\\
$R,g$ & Acceleration factor and geometry factor & Dimensionless\\
$\vect\theta$ & Estimated parameter vector & Component-dependent\\
$\chi,\delta$ & Susceptibility; chemical shift & Dimensionless; ppm if stated\\
\bottomrule
\end{longtable}
\renewcommand{\arraystretch}{1}
\index{Fourier transform!convention}
\index{gyromagnetic ratio}
\index{rotating frame}
\index{extended phase graph!convention}

\mainmatter

\part{Mathematical Foundations and Spin Dynamics}
% Scope: Develop only the mathematics and electromagnetism needed for spin dynamics, encoding, noise and estimation; derive prerequisites at their first useful level.
% Status: architecture only; no chapter prose or completed problems yet.
% Prerequisites: None; conventions in the front matter.
% Planned worked examples: Phasor signs, sampling limits and covariance propagation.
% Planned figures: Complex plane, convolution, sampling replicas.
% Planned challenge problems: Prove an aliasing ambiguity; optimize a noisy linear estimator.
\chapter{Mathematical and Physical Foundations}
\label{ch:01}

\section{Complex representation and dimensional reasoning}
\label{sec:01:complex-representation-and-dimensional-reasoning}
\subsection{Complex numbers, Euler's formula and phasors}
\subsection{Frequency, angular frequency and phase}
\subsection{SI units, scaling and dimensional checks}

\section{Linear algebra for spin and encoding}
\label{sec:01:linear-algebra-for-spin-and-encoding}
\subsection{Vector spaces, inner products and adjoints}
\subsection{Eigenvalues, eigenvectors and matrix exponentials}
\subsection{Singular values, rank and conditioning}

\section{Fourier analysis and linear systems}
\label{sec:01:fourier-analysis-and-linear-systems}
\subsection{Transform pairs and distributions}
\subsection{Convolution, filtering and impulse responses}
\subsection{Sampling, Nyquist limits and discrete transforms}

\section{Probability and measurement noise}
\label{sec:01:probability-and-measurement-noise}
\subsection{Random variables, expectations and covariance}
\subsection{Gaussian complex noise and magnitude statistics}
\subsection{Likelihood, bias and propagation of uncertainty}

\section{Estimation and experimental information}
\label{sec:01:estimation-and-experimental-information}
\subsection{Least squares and maximum likelihood}
\subsection{Fisher information and the Cramer--Rao bound}
\subsection{Identifiability and nuisance parameters}

\section{Fields, dipoles and angular momentum}
\label{sec:01:fields-dipoles-and-angular-momentum}
\subsection{Maxwell equations and quasistatic limits}
\subsection{Magnetic dipole energy and torque}
\subsection{Angular momentum and rotating coordinates}

\section{Worked examples}
\label{sec:01:worked-examples}
% Planned calculations: Phasor signs, sampling limits and covariance propagation.
\subsection{Derivation and quantitative calculation}
\subsection{Assumptions, limiting cases and scanner interpretation}

\section{Challenge problems}
\label{sec:01:challenge-problems}
% Problem design: Prove an aliasing ambiguity; optimize a noisy linear estimator.
% Use challengeproblem and stable labels prob:01:<topic>.
% Complete solutions belong in Appendix E, section for this chapter.
% Use \begin{solution}{prob:01:<topic>} ... \end{solution} there.
\subsection{Derivation and quantitative design}
\subsection{Model criticism and integrated reasoning}

% Scope: Spin-half quantum mechanics, thermal ensembles and the classical connection.
% Status: foundational chapter with explicit density-matrix dynamics.
% Prerequisites: Chapter 1 and the front-matter sign conventions.
% Planned figures: Implemented Zeeman levels and Bloch-vector geometry.
% Planned challenge problems: Implemented; solutions are in Appendix E.
\chapter{Quantum Spin Physics}
\label{ch:02}

Nuclear spin is intrinsic angular momentum, not a charged sphere whose
mechanical rotation must be inferred. Its state space, observables and
dynamical law determine what an MRI experiment can prepare and measure.
For uncoupled spin-half nuclei these rules admit an exact three-component
representation of the density operator. This explains why a classical-looking
magnetization vector describes many MRI experiments, even though thermal
polarization and the magnetic moment are quantum mechanical. The operator
framework used here is standard NMR theory \cite{levitt2008}.

\section{Spin as intrinsic angular momentum}
\label{sec:02:spin-as-intrinsic-angular-momentum}
\subsection{States, measurement bases and superposition}
Choose eigenvectors of $\op S_z$ as an orthonormal basis:
$\ket{\uparrow}=(1,0)^{\mathsf T}$ and
$\ket{\downarrow}=(0,1)^{\mathsf T}$.
A pure state is $\ket\psi=a\ket{\uparrow}+b\ket{\downarrow}$,
where $\abs a^2+\abs b^2=1$. A common phase multiplying both coefficients
changes no measurement probability. Their relative phase does:
transverse-basis measurements depend on $a^*b$.

Up to a common phase,
\[
 \ket\psi=\cos(\vartheta/2)\ket{\uparrow}
       +\ee^{\ii\varphi}\sin(\vartheta/2)\ket{\downarrow}.
\]
A $z$ measurement gives $+\hbar/2$ with probability
$\cos^2(\vartheta/2)$ and $-\hbar/2$ with probability
$\sin^2(\vartheta/2)$. Neither outcome implies a pre-existing classical
vector with simultaneous definite Cartesian components. Repeated
measurements on identically prepared systems determine expectation
values, which can be represented by a vector.

\subsection{Pauli matrices and spin operators}
The Pauli matrices and dimensional spin operators are
\begin{equation}
 \sigma_x=\begin{pmatrix}0&1\\1&0\end{pmatrix},\quad
 \sigma_y=\begin{pmatrix}0&-\ii\\\ii&0\end{pmatrix},\quad
 \sigma_z=\begin{pmatrix}1&0\\0&-1\end{pmatrix},\qquad
 \op S_j=\frac{\hbar}{2}\sigma_j .
 \label{eq:02:pauli}
\end{equation}
Direct multiplication gives
$\sigma_j\sigma_k=\delta_{jk}\mat I+
\ii\sum_l\epsilon_{jkl}\sigma_l$ and hence
$\op{\vect S}^{\,2}=3\hbar^2\mat I/4$.
An individual component eigenvalue has magnitude $\hbar/2$;
the squared total angular momentum has eigenvalue
$\hbar^2s(s+1)$ with $s=1/2$. Confusing these magnitudes gives an
incorrect magnetic moment.

For any unit vector $\vect n$,
$(\vect n\cdot\vect\sigma)^2=\mat I$, so the two projectors are
$\mat P_\pm=(\mat I\pm\vect n\cdot\vect\sigma)/2$.
Their expectation values supply measurement probabilities along an
arbitrary axis without explicitly constructing its eigenvectors.

\subsection{Commutators, uncertainty and eigenstates}
The multiplication law implies
$[\op S_j,\op S_k]=\ii\hbar\sum_l\epsilon_{jkl}\op S_l$.
Positivity of the norm of a linear combination of
$(\op A-\langle A\rangle)\ket\psi$ and
$(\op B-\langle B\rangle)\ket\psi$ gives
$\Delta A\,\Delta B\geq\tfrac12\abs{\langle[\op A,\op B]\rangle}$.
In $\ket{\uparrow}$, $\Delta S_x=\Delta S_y=\hbar/2$:
their product saturates the bound $\hbar^2/4$.
Even a fully polarized spin has transverse measurement uncertainty.

The states $\ket{\pm x}=(\ket{\uparrow}\pm\ket{\downarrow})/\sqrt2$
have equal $z$ populations but opposite transverse expectation values.
Similarly, $\ket{+y}=(\ket{\uparrow}+\ii\ket{\downarrow})/\sqrt2$.
Equal longitudinal populations do not imply zero transverse moment.

\section{Magnetic moments and Zeeman splitting}
\label{sec:02:magnetic-moments-and-zeeman-splitting}
\subsection{Gyromagnetic ratio and Hamiltonian}
The magnetic moment operator is $\op{\vect\mu}=\gamma\op{\vect S}$.
In a classical applied field,
\begin{equation}
 \op H(t)=-\gamma\op{\vect S}\cdot\vect B(t)
         =-\frac{\hbar\gamma}{2}\vect B(t)\cdot\vect\sigma .
 \label{eq:02:zeeman}
\end{equation}
The nuclear structure is incorporated in $\gamma$; it is not determined
by modeling the proton as a classical point charge. The units are energy
since $\hbar\gamma B$ is an action times inverse time.

\subsection{Energy levels and transition frequencies}
For positive $\gamma$ and field $B_0\vect e_z$,
\[
 E_\uparrow=-\hbar\omega_0/2,\qquad E_\downarrow=+\hbar\omega_0/2,
 \qquad \Delta E=\hbar\omega_0=hf_0 .
\]
A moment along the field has lower energy. A transverse perturbation
has off-diagonal matrix elements in this basis and drives transitions;
a longitudinal perturbation changes the splitting. Resonance matches
the drive to the splitting, including shielding and local-field shifts.

\begin{figure}[tb]
 \centering\begin{tikzpicture}[x=1cm,y=1cm,font=\small]
\draw[mri axis] (0,0)--(0,3.8) node[above] {$E$};
\draw[mri signal] (.4,.7)--(3,.7);
\draw[mri signal] (.4,3)--(3,3);
\node[anchor=west] at (.5,.35) {$\ket{\uparrow},\ -\hbar\omega_0/2$};
\node[anchor=west] at (.5,3.35) {$\ket{\downarrow},\ +\hbar\omega_0/2$};
\draw[<->,thick] (1.7,.8)--(1.7,2.9) node[midway,right] {RF: $\hbar\omega_0$};
\begin{scope}[xshift=7cm,yshift=1.7cm]
\draw[gray] (0,0) circle (1.55);
\draw[mri guide] (0,0) ellipse (1.55 and .42);
\draw[mri axis] (0,-1.7)--(0,1.9) node[above] {$p_z$};
\draw[mri axis] (-1.8,0)--(1.95,0) node[right] {$p_x$};
\draw[mri axis] (.8,.7)--(-1.3,-1.1) node[below] {$p_y$};
\draw[mri rf,-{Stealth}] (0,0)--(.85,1.15) node[right] {$\vect p$};
\node at (0,-2) {$\norm{\vect p}\leq1$};
\end{scope}
\end{tikzpicture}

 \caption{Positive-$\gamma$ Zeeman levels and a spin-half Bloch vector.
 RF couples the levels. The vector represents expectation values, not
 simultaneous sharp Cartesian spin components.}
 \label{fig:02:zeeman}
\end{figure}

\subsection{Positive and negative gyromagnetic ratios}
For negative $\gamma$, the ordering of the $S_z$ energies and the
precession handedness both reverse. The transition-frequency magnitude
is $\abs\gamma B_0/(2\pi)$. Equilibrium magnetization still points along
the field in this independent-spin paramagnetic model: the preferred
spin orientation and its conversion to magnetic moment both reverse.
A susceptibility proportional to $\gamma^2$ remains positive.
Our proton RF conventions require translation for negative-$\gamma$
nuclei.

\section{Unitary time evolution}
\label{sec:02:unitary-time-evolution}
\subsection{Schrodinger evolution and propagators}
Schrodinger evolution is
$\ii\hbar\partial_t\ket\psi=\op H(t)\ket\psi$.
For constant $\op H$, the solution is
$U(t)\ket{\psi(0)}$ with $U=\exp(-\ii\op Ht/\hbar)$.
Thus
\begin{equation}
 U_0(t)=\begin{pmatrix}
 \ee^{+\ii\omega_0t/2}&0\\0&\ee^{-\ii\omega_0t/2}
 \end{pmatrix}.
 \label{eq:02:free-u}
\end{equation}
The signs follow from the negative magnetic-dipole Hamiltonian.
A $2\pi$ rotation changes a spinor's sign, but returns its density matrix
and single-spin observables to their original values. It does not reverse
the MRI magnetization after a $360^\circ$ RF rotation.

\subsection{Relative phase and Larmor precession}
For amplitudes $a,b$,
$\langle S_x\rangle=\hbar\operatorname{Re}(a^*b)$,
$\langle S_y\rangle=\hbar\operatorname{Im}(a^*b)$ and
$\langle S_z\rangle=\hbar(\abs a^2-\abs b^2)/2$.
Since $a^*(t)b(t)=a^*(0)b(0)\ee^{-\ii\omega_0t}$,
\begin{equation}
 \langle S_x+\ii S_y\rangle(t)
 =\langle S_x+\ii S_y\rangle(0)\ee^{-\ii\omega_0t}.
 \label{eq:02:coherence}
\end{equation}
An initial $\ket{+x}$ state moves toward $-y$.
This is an independent physical check on the precession sign.

\subsection{Expectation values and the classical torque equation}
Use $\dd\langle\op S_i\rangle/\dd t
=(\ii/\hbar)\langle[\op H,\op S_i]\rangle$ and the commutators to obtain
\begin{equation}
 \frac{\dd}{\dd t}\langle\op{\vect S}\rangle
       =\gamma\langle\op{\vect S}\rangle\times\vect B .
 \label{eq:02:ehrenfest}
\end{equation}
Closure is exact because the Hamiltonian is linear in spin.
Multiplication by spin density and $\gamma$ gives the torque part of
the Bloch equations. A coupled-spin term such as
$\op S_{1z}\op S_{2z}$ instead makes first-moment evolution depend on
correlations. A single three-component vector then need not close.

\section{Density operators and ensembles}
\label{sec:02:density-operators-and-ensembles}
\subsection{Pure states, mixtures and coherence}
For a mixture with classical probabilities $p_a$,
$\op\rho=\sum_a p_a\ket{\psi_a}\bra{\psi_a}$.
It is Hermitian, positive semidefinite and trace one.
The mean of $\op A$ is $\Tr(\op\rho\op A)$.
Different ensemble decompositions can give the same density matrix;
single-spin measurements cannot distinguish them.

A pure state has $\op\rho^2=\op\rho$ and $\Tr(\op\rho^2)=1$.
An equal incoherent up/down mixture is $\mat I/2$, whereas
$\ket{+x}\bra{+x}$ has off-diagonal entries $1/2$.
Both have equal $z$ populations. Off-diagonal coherence is basis-dependent;
it is not an extra population of some third spin species.

\subsection{Bloch-vector representation and positivity}
Every Hermitian trace-one $2\times2$ matrix can be written
\begin{equation}
 \op\rho=\frac12(\mat I+\vect p\cdot\vect\sigma)
 =\frac12\begin{pmatrix}1+p_z&p_x-\ii p_y\\p_x+\ii p_y&1-p_z\end{pmatrix}.
 \label{eq:02:density-bloch}
\end{equation}
Here $p_j=\Tr(\op\rho\sigma_j)$ is dimensionless polarization and
$\langle\op{\vect S}\rangle=\hbar\vect p/2$.
The eigenvalues are $(1\pm\norm{\vect p})/2$, using
$(\vect p\cdot\vect\sigma)^2=\norm{\vect p}^2\mat I$.
Physicality requires $\norm{\vect p}\leq1$.
The Bloch sphere surface is pure; its interior is mixed, with purity
$(1+\norm{\vect p}^2)/2$.

For local spin density $\rho_s$,
\begin{equation}
 \vect M=\rho_s\gamma\frac{\hbar}{2}\vect p,\qquad
 M_+=\rho_s\gamma\hbar\,\rho_{21}.
 \label{eq:02:magnetization}
\end{equation}
The \emph{lower} off-diagonal matrix entry encodes $M_+$ in our convention.
These relations apply to a representative local one-spin density matrix,
including an ensemble-averaged reduced state.

\subsection{Dephasing versus loss of single-spin coherence}
Unitary evolution obeys
$\dot{\op\rho}=-(\ii/\hbar)[\op H,\op\rho]$ and preserves eigenvalues
and purity. Frequency dispersion can nevertheless reduce ensemble
transverse mean:
$p_+(t)=p_+(0)\int P(\Delta\omega)\ee^{-\ii\Delta\omega t}\dd\Delta\omega$.
Each member retains magnitude while their vector sum shrinks.
Offset knowledge or a refocusing pulse can undo this cancellation.

Environmental coupling can reduce coherence in a spin's reduced density
matrix too, and this need not be undone by a simple pulse.
Some slowly varying environmental perturbations are refocusable.
The relevant distinctions are correlation time and controllability;
dephasing is not synonymous with an irrecoverable process in every model.

\section{Thermal equilibrium and polarization}
\label{sec:02:thermal-equilibrium-and-polarization}
\subsection{Partition function and Boltzmann populations}
For independent spins at bath temperature $T_{\rm bath}$, let
$\beta=1/(k_{\rm B}T_{\rm bath})$ and $x=\beta\hbar\omega_0/2$.
Then
\[
 Z=\Tr(\ee^{-\beta\op H_0})=2\cosh x,\qquad
 \op\rho_{\rm eq}=\frac{\ee^{x\sigma_z}}{2\cosh x}
                =\frac12(\mat I+\tanh x\,\sigma_z).
\]
The probabilities are $\ee^{\pm x}/(2\cosh x)$, giving
\begin{equation}
 p_\uparrow-p_\downarrow=\tanh x,\qquad
 p_\uparrow/p_\downarrow=\ee^{2x}.
 \label{eq:02:boltzmann}
\end{equation}
The population difference divided by total population differs from the
fractional excess relative to one population. Their small-$x$ limits
differ by a factor of two.

\subsection{High-temperature expansion and its accuracy}
For $\abs x\ll1$, $\tanh x=x-x^3/3+O(x^5)$.
Replacing it by $x$ has relative error approximately $x^2/3$.
Physiological temperature is high relative to the proton Zeeman splitting,
not necessarily to every microscopic energy in tissue.
The identity part of the thermal density matrix greatly exceeds the
polarized deviation. Unitary RF evolution leaves the identity unchanged;
it rotates only that small deviation.

\subsection{Spin density and macroscopic magnetization}
Combining the population and moment expressions gives
\begin{equation}
 M_0=\rho_s\frac{\gamma\hbar}{2}
       \tanh\left(\frac{\hbar\gamma B_0}{2k_{\rm B}T_{\rm bath}}\right)
 \simeq \frac{\rho_s\gamma^2\hbar^2B_0}{4k_{\rm B}T_{\rm bath}} .
 \label{eq:02:curie}
\end{equation}
Units are moment per volume, or A/m. The linear field scaling assumes
independent spins and small polarization. Received voltage also depends
on frequency and coil response; SNR does not follow from $M_0$ alone.

For $N$ independent spins, mean population imbalance is $N\tanh x$
and its standard deviation is $\sqrt{N(1-\tanh^2x)}$.
Macroscopic MRI volumes therefore have large absolute net polarization
despite its small fraction. This is distinct from the question of whether
receiver noise is dominated by spin fluctuations or electrical noise.

\section{Driven spins and the classical limit}
\label{sec:02:driven-spins-and-the-classical-limit}
\subsection{RF interaction Hamiltonian}
For a co-rotating field of amplitude $B_a$ and phase $\phi$, the
rotating-frame Hamiltonian under the rotating-wave approximation is
\[
 \op H_{\rm rot}=-\frac{\hbar}{2}\left[
 \Delta\omega\sigma_z+\omega_1(\cos\phi\,\sigma_x+\sin\phi\,\sigma_y)\right],
 \qquad\omega_1=\gamma B_a .
\]
\Cref{ch:04} derives the transformation and $B_a=\abs{\Btx}$.
On resonance, for constant phase, the propagator is
$U=\cos(\alpha/2)\mat I+\ii\sin(\alpha/2)\vect n\cdot\vect\sigma$,
with $\alpha=\omega_1t$ and $\vect n=(\cos\phi,\sin\phi,0)$.
Separate even and odd exponential powers to derive this formula.

An initially spin-up pure state then has spin-down probability
$\sin^2(\alpha/2)$. A thermal state instead has its polarization rotated:
for $\phi=0$,
$\vect p=(0,p_0\sin\alpha,p_0\cos\alpha)$.
A $90^\circ$ pulse equalizes populations while maximizing transverse
magnetization; population imbalance alone no longer characterizes it.

\subsection{Rotating-wave approximation and resonance}
A linear RF oscillation has equal co- and counter-rotating components.
In a resonant frame one is slow; the other oscillates near twice the carrier.
When drive strength, detuning and envelope variation rates are small
relative to the carrier, the latter averages to a small correction.
A linearly polarized peak field $B_{\rm pk}$ therefore gives
$\omega_1=\gamma B_{\rm pk}/2$. For an ideal purely co-rotating field,
discarding a counter-rotating term is unnecessary. Strong drives or
rapid modulation require the full Hamiltonian.

\subsection{Ensemble Bloch dynamics and hidden microscopic physics}
For a single spin with a linear Hamiltonian the Bloch vector exactly
represents the density matrix, including mixed states.
Relaxation requires an open-system model. In a Markovian spin-half
model with energy-downward and energy-upward rates $W_d,W_u$ and
pure-dephasing rate $\Gamma_\phi$,
\[
 R_1=W_d+W_u,\qquad R_2=R_1/2+\Gamma_\phi .
\]
For protons the lower-energy state is spin up. Detailed balance fixes
the equilibrium polarization. The bound $T_2\leq2T_1$ follows for this
model; empirical effective times in heterogeneous multi-pool data need
not describe this same pair of microscopic processes.

Classical ensemble magnetization suffices when first moments close under
the dynamics. It hides quantum measurement uncertainty, purity and
higher correlations. Scalar couplings in spectroscopy can transfer
information into correlations outside a one-vector model. Routine imaging
often admits local Bloch vectors with relaxation and transport as an
appropriate reduced description.

\section{Worked examples}
\label{sec:02:worked-examples}
\subsection{Derivation and quantitative calculation}
\begin{workedexample}[Thermal polarization]
Use $\gammabar=\SI{42.57748}{\mega\hertz\per\tesla}$,
$\hbar=\SI{1.054571817e-34}{\joule\second}$ and
$k_{\rm B}=\SI{1.380649e-23}{\joule\per\kelvin}$.
At $T_{\rm bath}=\SI{310}{\kelvin}$:
\begin{center}
\begin{tabular}{@{}rrr@{}}\toprule
$B_0$ (T)&$f_0$ (MHz)&$10^6(p_\uparrow-p_\downarrow)$\\\midrule
1.5&63.8662&4.9437\\3.0&127.7324&9.8874\\7.0&298.0424&23.0706\\\bottomrule
\end{tabular}
\end{center}
The rounded free-proton value illustrates scale, not a scanner calibration;
shielding slightly changes actual resonance.
At 7 T the relative high-temperature error is only
$x^2/3\simeq1.77\times10^{-10}$.
\end{workedexample}
\begin{workedexample}[Net moment in a voxel]
For illustrative water-like proton density
$\rho_s=\SI{6.69e28}{\per\cubic\metre}$ at 3 T and 310 K,
$M_0=\SI{9.33e-3}{\ampere\per\metre}$.
A \SI{1}{\cubic\milli\metre} voxel has $6.69\times10^{19}$ protons,
mean population excess $6.62\times10^{14}$ and equilibrium magnetic
moment $\SI{9.33e-12}{\ampere\square\metre}$.
The density is a stated model input, not a universal tissue constant.
\end{workedexample}
\subsection{Assumptions, limiting cases and scanner interpretation}
\begin{workedexample}[Equal populations with coherence]
Apply $U=\ee^{+\ii\pi\sigma_x/4}$ to
$\op\rho_{\rm eq}=(\mat I+p_0\sigma_z)/2$.
Multiplication gives $\op\rho'=(\mat I+p_0\sigma_y)/2$:
both diagonal entries are $1/2$, while $\rho'_{21}=\ii p_0/2$.
Thus $M_+=\ii M_0$. Replacing this state by an incoherent 50:50
mixture would incorrectly predict no signal.
\end{workedexample}

\section{Challenge problems}
\label{sec:02:challenge-problems}
\subsection{Derivation and quantitative design}
\begin{challengeproblem}[A proposed density matrix]
\label{prob:02:tomography}
For $\op\rho=\tfrac12\begin{pmatrix}1+a&b-\ii c\\b+\ii c&1-a\end{pmatrix}$,
derive physicality, purity and Cartesian-axis measurement probabilities.
Evaluate $(a,b,c)=(0.6,0.3,0.4)$ and explain the failure of
$(0.9,0.6,0)$. Distinguish a nonphysical fit from a negative noisy datum.
\end{challengeproblem}
\begin{challengeproblem}[Population conventions]
\label{prob:02:thermal}
Derive the exact spin-half partition function, polarization and magnetization
for signed $\gamma$. Show that equilibrium magnetization is parallel to
$B_0$ for either sign. At 3 T and 310 K compare
$(N_\uparrow-N_\downarrow)/(N_\uparrow+N_\downarrow)$ with
$(N_\uparrow-N_\downarrow)/N_\downarrow$ for protons.
Estimate when linearized polarization has one-percent relative error.
\end{challengeproblem}
\begin{challengeproblem}[A two-pulse interferometer]
\label{prob:02:ramsey}
Starting from $p_0\vect e_z$, apply a phase-zero $\pi/2$ pulse,
free rotating-frame evolution for $\tau$ with detuning $\Delta\omega$,
then a $\pi/2$ pulse of phase $\phi$. Neglect relaxation.
Derive final $p_z$ by quantum propagators and by Bloch rotations.
\end{challengeproblem}
\subsection{Model criticism and integrated reasoning}
\begin{challengeproblem}[Reversible ensemble dephasing]
\label{prob:02:dephase}
Equal subensembles start in $\ket{+x}$ with offsets $+\Omega$ and
$-\Omega$. Find their average density matrix and purity versus time.
Apply a phase-zero $\pi$ pulse at $\tau$ and determine the state at
$2\tau$. Can loss of ensemble purity prove irreversible decoherence?
\end{challengeproblem}
\begin{challengeproblem}[A physical relaxation map]
\label{prob:02:rates}
For energy-downward and energy-upward rates $W_d,W_u$, derive equilibrium
polarization and $T_1$ from the population master equation.
If coherence decays at $(W_d+W_u)/2+\Gamma_\phi$, find $T_2$ and
the allowed ratio $T_2/T_1$. Why does fitting two arbitrary exponentials
not identify a valid microscopic model automatically?
\end{challengeproblem}

% Scope: Construct the phenomenological Bloch model from torque and equilibration; solve free evolution and connect reversibility, relaxation and exchange.
% Status: architecture only; no chapter prose or completed problems yet.
% Prerequisites: Chapters 1 and 2.
% Planned worked examples: Longitudinal recovery, transverse decay and an inhomogeneous ensemble.
% Planned figures: Relaxation curves and isochromat refocusing.
% Planned challenge problems: Solve a piecewise field history; distinguish static dephasing from irreversible loss.
\chapter{Classical Magnetization and the Bloch Equations}
\label{ch:03}

\section{From magnetic torque to ensemble dynamics}
\label{sec:03:from-magnetic-torque-to-ensemble-dynamics}
\subsection{Torque equation and conserved magnitude}
\subsection{Equilibrium magnetization and phenomenological relaxation}
\subsection{Assumptions of the single-pool Bloch model}

\section{Solutions in a static field}
\label{sec:03:solutions-in-a-static-field}
\subsection{Longitudinal recovery and inversion}
\subsection{Complex transverse solution and phase}
\subsection{Affine propagators and piecewise-constant fields}

\section{Inhomogeneous ensembles}
\label{sec:03:inhomogeneous-ensembles}
\subsection{Frequency distributions and free induction decay}
\subsection{T2, T2-star and reversible broadening}
\subsection{Limits of exponential effective relaxation}

\section{Driven and time-dependent dynamics}
\label{sec:03:driven-and-time-dependent-dynamics}
\subsection{Constant effective-field solution}
\subsection{Numerical integration and operator splitting}
\subsection{Accuracy, stability and conservation checks}

\section{Beyond independent stationary spins}
\label{sec:03:beyond-independent-stationary-spins}
\subsection{Bloch--Torrey transport and diffusion}
\subsection{Bloch--McConnell exchange structure}
\subsection{Where phenomenology requires microscopic input}

\section{Worked examples}
\label{sec:03:worked-examples}
% Planned calculations: Longitudinal recovery, transverse decay and an inhomogeneous ensemble.
\subsection{Derivation and quantitative calculation}
\subsection{Assumptions, limiting cases and scanner interpretation}

\section{Challenge problems}
\label{sec:03:challenge-problems}
% Problem design: Solve a piecewise field history; distinguish static dephasing from irreversible loss.
% Use challengeproblem and stable labels prob:03:<topic>.
% Complete solutions belong in Appendix E, section for this chapter.
% Use \begin{solution}{prob:03:<topic>} ... \end{solution} there.
\subsection{Derivation and quantitative design}
\subsection{Model criticism and integrated reasoning}

% Scope: Establish exact frame transformations and RF phase conventions before treating excitation, inversion and selective pulses.
% Status: architecture only; no chapter prose or completed problems yet.
% Prerequisites: Chapters 2 and 3.
% Planned worked examples: Rotating-frame transformation and flip angle of a rectangular RF pulse.
% Planned figures: Laboratory and rotating axes; effective-field trajectories.
% Planned challenge problems: Derive a detuned pulse propagator; debug a factor-of-two RF amplitude error.
\chapter{RF Excitation and the Rotating Frame}
\label{ch:04}

\section{Rotating coordinate systems}
\label{sec:04:rotating-coordinate-systems}
\subsection{Active rotations versus passive coordinates}
\subsection{Transformation of time derivatives}
\subsection{Detuning and the effective field}

\section{Circular RF components}
\label{sec:04:circular-rf-components}
\subsection{Real fields and complex phasors}
\subsection{Co-rotating B1-plus and counter-rotating fields}
\subsection{Linear polarization and the rotating-wave approximation}

\section{Hard-pulse dynamics}
\label{sec:04:hard-pulse-dynamics}
\subsection{Flip angle, RF phase and rotation axis}
\subsection{On-resonance excitation and inversion}
\subsection{Off-resonance excitation and bandwidth}

\section{Pulse imperfections and calibration}
\label{sec:04:pulse-imperfections-and-calibration}
\subsection{Transmit inhomogeneity and flip-angle errors}
\subsection{Finite-duration effects and relaxation}
\subsection{Phase cycling and receiver phase}

\section{Selective excitation as an introduction}
\label{sec:04:selective-excitation-as-an-introduction}
\subsection{Slice-selection gradients and RF bandwidth}
\subsection{Small-tip response and pulse envelopes}
\subsection{Boundaries of the hard-pulse approximation}

\section{Worked examples}
\label{sec:04:worked-examples}
% Planned calculations: Rotating-frame transformation and flip angle of a rectangular RF pulse.
\subsection{Derivation and quantitative calculation}
\subsection{Assumptions, limiting cases and scanner interpretation}

\section{Challenge problems}
\label{sec:04:challenge-problems}
% Problem design: Derive a detuned pulse propagator; debug a factor-of-two RF amplitude error.
% Use challengeproblem and stable labels prob:04:<topic>.
% Complete solutions belong in Appendix E, section for this chapter.
% Use \begin{solution}{prob:04:<topic>} ... \end{solution} there.
\subsection{Derivation and quantitative design}
\subsection{Model criticism and integrated reasoning}


\part{Encoding, Signal Formation and Pulse Sequences}
% Scope: Derive spatial phase accumulation and trajectory geometry, with sampling definitions separated from effective resolution.
% Status: architecture only; no chapter prose or completed problems yet.
% Prerequisites: Chapters 1, 3 and 4.
% Planned worked examples: Slice thickness, gradient strength, FOV, k-space extent and pixel spacing.
% Planned figures: Slice selection, Cartesian, radial, spiral and EPI trajectories.
% Planned challenge problems: Design a constrained Cartesian readout; expose a zero-filling resolution fallacy.
\chapter{Spatial Encoding and k-Space}
\label{ch:05}

\section{Gradient fields and spatial phase}
\label{sec:05:gradient-fields-and-spatial-phase}
\subsection{Linear field expansion and gradient units}
\subsection{Phase accumulation and trajectory origin}
\subsection{Refocusing pulses and effective gradient signs}

\section{Slice selection and partition encoding}
\label{sec:05:slice-selection-and-partition-encoding}
\subsection{Frequency-to-position mapping}
\subsection{Slice rephasing and finite pulse duration}
\subsection{Two-dimensional slices versus three-dimensional slabs}

\section{Frequency and phase encoding}
\label{sec:05:frequency-and-phase-encoding}
\subsection{Readout gradients and dwell time}
\subsection{Phase-encoding steps and Cartesian grids}
\subsection{Echo position and prephasing}

\section{Sampling geometry}
\label{sec:05:sampling-geometry}
\subsection{Field of view and k-space spacing}
\subsection{Finite extent, nominal resolution and voxel grid}
\subsection{Nyquist sampling, aliasing and zero filling}

\section{Trajectory families and ordering}
\label{sec:05:trajectory-families-and-ordering}
\subsection{Cartesian and echo-planar trajectories}
\subsection{Radial and spiral sampling}
\subsection{Linear, centric and interleaved view ordering}

\section{Gradient moments and trajectory fidelity}
\label{sec:05:gradient-moments-and-trajectory-fidelity}
\subsection{Zeroth and higher moments}
\subsection{Gradient delays and eddy-current errors}
\subsection{Concomitant fields and nonlinear encoding}

\section{Worked examples}
\label{sec:05:worked-examples}
% Planned calculations: Slice thickness, gradient strength, FOV, k-space extent and pixel spacing.
\subsection{Derivation and quantitative calculation}
\subsection{Assumptions, limiting cases and scanner interpretation}

\section{Challenge problems}
\label{sec:05:challenge-problems}
% Problem design: Design a constrained Cartesian readout; expose a zero-filling resolution fallacy.
% Use challengeproblem and stable labels prob:05:<topic>.
% Complete solutions belong in Appendix E, section for this chapter.
% Use \begin{solution}{prob:05:<topic>} ... \end{solution} there.
\subsection{Derivation and quantitative design}
\subsection{Model criticism and integrated reasoning}

% Scope: Derive reception and the MRI forward model, then identify exactly when Fourier inversion is justified; establish PSF and complex-noise language.
% Status: architecture only; no chapter prose or completed problems yet.
% Prerequisites: Chapters 1 through 5.
% Planned worked examples: Two-point object, finite readout blur and noise through a DFT.
% Planned figures: Receive sensitivity, Fourier pairs and point-spread functions.
% Planned challenge problems: Derive the forward operator for a decaying readout; compare sampling and resolution.
\chapter{Signal Formation and Fourier Imaging}
\label{ch:06}

\section{From precession to received voltage}
\label{sec:06:from-precession-to-received-voltage}
\subsection{Faraday induction and reciprocity}
\subsection{Complex demodulation and receiver phase}
\subsection{Receive sensitivity and signal normalization}

\section{The MRI signal equation}
\label{sec:06:the-mri-signal-equation}
\subsection{Spatial integration of transverse magnetization}
\subsection{Encoding, relaxation and off-resonance during readout}
\subsection{Motion, multiple species and model limitations}

\section{The Fourier imaging limit}
\label{sec:06:the-fourier-imaging-limit}
\subsection{Stationary object and time-independent weighting}
\subsection{Forward and inverse transform conventions}
\subsection{Discrete sampling and voxel integration}

\section{Image formation and point-spread functions}
\label{sec:06:image-formation-and-point-spread-functions}
\subsection{Finite support and truncation}
\subsection{Apodization, convolution and effective resolution}
\subsection{Center and periphery without oversimplification}

\section{Incomplete and modified Fourier data}
\label{sec:06:incomplete-and-modified-fourier-data}
\subsection{Partial Fourier and conjugate-symmetry assumptions}
\subsection{Phase estimation and homodyne concepts}
\subsection{Zero padding, interpolation and display grids}

\section{Complex images and noise}
\label{sec:06:complex-images-and-noise}
\subsection{DFT normalization and covariance}
\subsection{Phase images, magnitude images and noise floors}
\subsection{Receive weighting and the meaning of image intensity}

\section{Worked examples}
\label{sec:06:worked-examples}
% Planned calculations: Two-point object, finite readout blur and noise through a DFT.
\subsection{Derivation and quantitative calculation}
\subsection{Assumptions, limiting cases and scanner interpretation}

\section{Challenge problems}
\label{sec:06:challenge-problems}
% Problem design: Derive the forward operator for a decaying readout; compare sampling and resolution.
% Use challengeproblem and stable labels prob:06:<topic>.
% Complete solutions belong in Appendix E, section for this chapter.
% Use \begin{solution}{prob:06:<topic>} ... \end{solution} there.
\subsection{Derivation and quantitative design}
\subsection{Model criticism and integrated reasoning}

% Scope: Build FID, spin echo, gradient echo, spoiled GRE and inversion recovery from explicit magnetization histories; establish the sequence vocabulary used later.
% Status: architecture only; no chapter prose or completed problems yet.
% Prerequisites: Chapters 3 through 6.
% Planned worked examples: Spin-echo signal, spoiled-GRE signal, Ernst angle and inversion null time.
% Planned figures: FID, Hahn echo, GRE and inversion-recovery timing diagrams.
% Planned challenge problems: Infer a sequence from timing; optimize flip angle with competing tissue contrasts.
\chapter{Basic Pulse Sequences}
\label{ch:07}

\section{Reading a sequence diagram}
\label{sec:07:reading-a-sequence-diagram}
\subsection{RF, gradient and ADC event tracks}
\subsection{TR, TE, TI and timing references}
\subsection{Two-dimensional and three-dimensional acquisition loops}

\section{Free induction and Hahn spin echo}
\label{sec:07:free-induction-and-hahn-spin-echo}
\subsection{FID and static phase dispersion}
\subsection{Refocusing and echo formation}
\subsection{What a spin echo does and does not reverse}

\section{Gradient echoes}
\label{sec:07:gradient-echoes}
\subsection{Gradient reversal and echo time}
\subsection{T2-star weighting and signal phase}
\subsection{Comparison with RF refocusing}

\section{Spoiled gradient echo}
\label{sec:07:spoiled-gradient-echo}
\subsection{Longitudinal recurrence and steady state}
\subsection{FLASH and SPGR families}
\subsection{Ernst angle and the limits of perfect spoiling}

\section{Inversion and saturation recovery}
\label{sec:07:inversion-and-saturation-recovery}
\subsection{Finite-TR signal equations}
\subsection{Null times and imperfect inversion}
\subsection{STIR, FLAIR and contrast-selective suppression}

\section{Sequence comparisons}
\label{sec:07:sequence-comparisons}
\subsection{Transient versus steady-state acquisition}
\subsection{SNR efficiency and scan-time accounting}
\subsection{Clinical examples and hidden model assumptions}

\section{Worked examples}
\label{sec:07:worked-examples}
% Planned calculations: Spin-echo signal, spoiled-GRE signal, Ernst angle and inversion null time.
\subsection{Derivation and quantitative calculation}
\subsection{Assumptions, limiting cases and scanner interpretation}

\section{Challenge problems}
\label{sec:07:challenge-problems}
% Problem design: Infer a sequence from timing; optimize flip angle with competing tissue contrasts.
% Use challengeproblem and stable labels prob:07:<topic>.
% Complete solutions belong in Appendix E, section for this chapter.
% Use \begin{solution}{prob:07:<topic>} ... \end{solution} there.
\subsection{Derivation and quantitative design}
\subsection{Model criticism and integrated reasoning}

% Scope: Develop echo trains, coherent steady states and prepared acquisitions; introduce pathway effects physically and defer their full algebra to Chapter 13.
% Status: architecture only; no chapter prose or completed problems yet.
% Prerequisites: Chapters 5 through 7; Chapter 13 for pathway calculations.
% Planned worked examples: bSSFP off-resonance response, effective TE and prepared GRE transients.
% Planned figures: TSE, bSSFP, EPI and MPRAGE timing.
% Planned challenge problems: Choose echo ordering for a target PSF; analyze startup and preparation bias.
\chapter{Advanced Pulse Sequences}
\label{ch:08}

\section{Fast and turbo spin echo}
\label{sec:08:fast-and-turbo-spin-echo}
\subsection{CPMG conditions, RARE and FSE/TSE}
\subsection{Stimulated echoes and variable refocusing angles}
\subsection{Echo-train modulation, effective TE and blurring}

\section{Balanced steady-state free precession}
\label{sec:08:balanced-steady-state-free-precession}
\subsection{Balanced gradient moments and coherent steady state}
\subsection{Frequency response, phase cycling and banding}
\subsection{Startup transients and cardiac cine}

\section{Echo-planar imaging}
\label{sec:08:echo-planar-imaging}
\subsection{Gradient-echo and spin-echo EPI}
\subsection{Echo spacing, segmentation and readout duration}
\subsection{Odd-even phase errors and distortion mechanisms}

\section{Multi-echo and hybrid acquisitions}
\label{sec:08:multi-echo-and-hybrid-acquisitions}
\subsection{Multi-echo GRE and spin-echo trains}
\subsection{GRASE and mixed relaxation weighting}
\subsection{Temporal ordering and contrast consistency}

\section{Prepared three-dimensional imaging}
\label{sec:08:prepared-three-dimensional-imaging}
\subsection{Magnetization preparation and readout perturbation}
\subsection{MPRAGE and inversion efficiency}
\subsection{MP2RAGE combination and residual dependencies}

\section{Chemical-species and multislice strategies}
\label{sec:08:chemical-species-and-multislice-strategies}
\subsection{Dixon phase evolution and echo placement}
\subsection{Simultaneous multislice acquisition}
\subsection{Preparation modules and links to MT, diffusion and ASL}

\section{Worked examples}
\label{sec:08:worked-examples}
% Planned calculations: bSSFP off-resonance response, effective TE and prepared GRE transients.
\subsection{Derivation and quantitative calculation}
\subsection{Assumptions, limiting cases and scanner interpretation}

\section{Challenge problems}
\label{sec:08:challenge-problems}
% Problem design: Choose echo ordering for a target PSF; analyze startup and preparation bias.
% Use challengeproblem and stable labels prob:08:<topic>.
% Complete solutions belong in Appendix E, section for this chapter.
% Use \begin{solution}{prob:08:<topic>} ... \end{solution} there.
\subsection{Derivation and quantitative design}
\subsection{Model criticism and integrated reasoning}


\part{Contrast, Quantification and Signal Pathways}
% Scope: Connect relaxation models, exchange and susceptibility to sequence-dependent tissue contrast without treating tissue constants as universal.
% Status: architecture only; no chapter prose or completed problems yet.
% Prerequisites: Chapters 2, 3, 7 and 8.
% Planned worked examples: Competing tissue contrast and multi-pool relaxation.
% Planned figures: Spectral densities, exchange and contrast curves.
% Planned challenge problems: Test a monoexponential assumption; compare nulling under imperfect inversion.
\chapter{Contrast and Relaxation}
\label{ch:09}

\section{Microscopic origins of relaxation}
\label{sec:09:microscopic-origins-of-relaxation}
\subsection{Fluctuating fields and correlation functions}
\subsection{Spectral density, T1 and T2}
\subsection{Field strength and molecular-motion regimes}

\section{Susceptibility and transverse dephasing}
\label{sec:09:susceptibility-and-transverse-dephasing}
\subsection{Microscopic and macroscopic field variation}
\subsection{Reversible and irreversible signal loss}
\subsection{Spin-echo versus gradient-echo contrast}

\section{Exchange and magnetization transfer}
\label{sec:09:exchange-and-magnetization-transfer}
\subsection{Two-pool exchange models}
\subsection{Bound-pool saturation and MT contrast}
\subsection{MTR, quantitative MT and model dependence}

\section{Contrast preparation}
\label{sec:09:contrast-preparation}
\subsection{Inversion, saturation and T2 preparation}
\subsection{Fat suppression and spectral selectivity}
\subsection{Spin locking and T1-rho concepts}

\section{Contrast agents and tissue models}
\label{sec:09:contrast-agents-and-tissue-models}
\subsection{Relaxivity, concentration and nonlinear effects}
\subsection{Compartmentation and exchange limitations}
\subsection{Tissue variability and sequence dependence}

\section{Contrast interpretation}
\label{sec:09:contrast-interpretation}
\subsection{Weighting versus parameter measurement}
\subsection{Clinical uses of T1, T2 and susceptibility}
\subsection{Failure of simple contrast heuristics}

\section{Worked examples}
\label{sec:09:worked-examples}
% Planned calculations: Competing tissue contrast and multi-pool relaxation.
\subsection{Derivation and quantitative calculation}
\subsection{Assumptions, limiting cases and scanner interpretation}

\section{Challenge problems}
\label{sec:09:challenge-problems}
% Problem design: Test a monoexponential assumption; compare nulling under imperfect inversion.
% Use challengeproblem and stable labels prob:09:<topic>.
% Complete solutions belong in Appendix E, section for this chapter.
% Use \begin{solution}{prob:09:<topic>} ... \end{solution} there.
\subsection{Derivation and quantitative design}
\subsection{Model criticism and integrated reasoning}

% Scope: Derive diffusion encoding from stochastic motion and gradient waveforms; progress from ADC to tensor and limited non-Gaussian models with clear identifiability boundaries.
% Status: architecture only; no chapter prose or completed problems yet.
% Prerequisites: Chapters 1, 5, 7 and 9.
% Planned worked examples: Stejskal--Tanner b-value, general b-matrix and tensor metrics.
% Planned figures: Brownian propagators, diffusion gradients and tensor ellipsoids.
% Planned challenge problems: Design equal-b unequal-time experiments; diagnose a biased tensor fit.
\chapter{Diffusion MRI}
\label{ch:10}

\section{Stochastic motion and transport}
\label{sec:10:stochastic-motion-and-transport}
\subsection{Brownian motion and the diffusion equation}
\subsection{Gaussian propagator and mean-square displacement}
\subsection{Characteristic functions and displacement encoding}

\section{Gradient-waveform encoding}
\label{sec:10:gradient-waveform-encoding}
\subsection{Stejskal--Tanner experiment and finite pulses}
\subsection{General waveform b-value and b-matrix}
\subsection{Imaging-gradient cross-terms and diffusion time}

\section{Apparent diffusion and tensor imaging}
\label{sec:10:apparent-diffusion-and-tensor-imaging}
\subsection{ADC and log-signal fitting}
\subsection{Diffusion tensor, eigenvectors and eigenvalues}
\subsection{Mean diffusivity, fractional anisotropy and invariants}

\section{Beyond Gaussian diffusion}
\label{sec:10:beyond-gaussian-diffusion}
\subsection{Restricted diffusion and time dependence}
\subsection{Diffusion kurtosis and cumulant limits}
\subsection{Multi-compartment models and IVIM}

\section{Acquisition and correction}
\label{sec:10:acquisition-and-correction}
\subsection{Spin-echo EPI and diffusion preparation}
\subsection{Motion, eddy currents and susceptibility distortion}
\subsection{Gradient nonlinearity, b-vector rotation and noise-floor bias}

\section{Interpretation and advanced encoding}
\label{sec:10:interpretation-and-advanced-encoding}
\subsection{Acute ischemia and oncologic examples}
\subsection{Tractography, crossing fibers and inference limits}
\subsection{Oscillating gradients and tensor-valued encoding}

\section{Worked examples}
\label{sec:10:worked-examples}
% Planned calculations: Stejskal--Tanner b-value, general b-matrix and tensor metrics.
\subsection{Derivation and quantitative calculation}
\subsection{Assumptions, limiting cases and scanner interpretation}

\section{Challenge problems}
\label{sec:10:challenge-problems}
% Problem design: Design equal-b unequal-time experiments; diagnose a biased tensor fit.
% Use challengeproblem and stable labels prob:10:<topic>.
% Complete solutions belong in Appendix E, section for this chapter.
% Use \begin{solution}{prob:10:<topic>} ... \end{solution} there.
\subsection{Derivation and quantitative design}
\subsection{Model criticism and integrated reasoning}

% Scope: Organize quantitative measurements around forward models, calibration, identifiability and uncertainty; provide balanced coverage across parameter families.
% Status: architecture only; no chapter prose or completed problems yet.
% Prerequisites: Chapters 1 and 6 through 10.
% Planned worked examples: T1/T2 fits, B0 phase unwrapping, uncertainty and protocol optimization.
% Planned figures: Signal families, likelihood contours and parameter covariance.
% Planned challenge problems: Design identifiable multiparametric sampling; quantify calibration-induced bias.
\chapter{Quantitative MRI}
\label{ch:11}

\section{Quantification as an inverse problem}
\label{sec:11:quantification-as-an-inverse-problem}
\subsection{Forward models, nuisance parameters and calibration}
\subsection{Likelihoods for complex and magnitude data}
\subsection{Bias, precision, accuracy and repeatability}

\section{Longitudinal relaxometry}
\label{sec:11:longitudinal-relaxometry}
\subsection{Inversion and saturation recovery}
\subsection{Variable flip angles and B1-plus correction}
\subsection{Look--Locker readout perturbation and mapping strategies}

\section{Transverse relaxometry and proton density}
\label{sec:11:transverse-relaxometry-and-proton-density}
\subsection{Multi-echo T2 and stimulated-echo contamination}
\subsection{T2-star, nonexponential decay and echo sampling}
\subsection{Proton density, receive bias and scaling ambiguity}

\section{Field and transmit mapping}
\label{sec:11:field-and-transmit-mapping}
\subsection{B0 phase differences and phase unwrapping}
\subsection{B1-plus mapping and flip-angle calibration}
\subsection{Spatial regularization and calibration transfer}

\section{Susceptibility, exchange and perfusion}
\label{sec:11:susceptibility-exchange-and-perfusion}
\subsection{QSM forward dipole model and inversion ambiguity}
\subsection{MTR, quantitative MT and exchange estimates}
\subsection{Perfusion parameters and ASL quantification links}

\section{Multiparametric estimation}
\label{sec:11:multiparametric-estimation}
\subsection{Joint models and parameter coupling}
\subsection{Fisher information, uncertainty and noise propagation}
\subsection{Multi-parametric mapping and synthetic MRI}

\section{Validation and experimental design}
\label{sec:11:validation-and-experimental-design}
\subsection{Ground truth, phantoms and repeat scans}
\subsection{Model mismatch and residual analysis}
\subsection{Sampling design under scan-time constraints}

\section{Worked examples}
\label{sec:11:worked-examples}
% Planned calculations: T1/T2 fits, B0 phase unwrapping, uncertainty and protocol optimization.
\subsection{Derivation and quantitative calculation}
\subsection{Assumptions, limiting cases and scanner interpretation}

\section{Challenge problems}
\label{sec:11:challenge-problems}
% Problem design: Design identifiable multiparametric sampling; quantify calibration-induced bias.
% Use challengeproblem and stable labels prob:11:<topic>.
% Complete solutions belong in Appendix E, section for this chapter.
% Use \begin{solution}{prob:11:<topic>} ... \end{solution} there.
\subsection{Derivation and quantitative design}
\subsection{Model criticism and integrated reasoning}

% Scope: Treat MRF as one transient quantitative acquisition family; connect spin dynamics, sampling and inference without duplicating the general qMRI or inverse-problem chapters.
% Status: architecture only; no chapter prose or completed problems yet.
% Prerequisites: Chapters 8 and 11; Chapters 13 and 16 for advanced implementations.
% Planned worked examples: Normalized matching, proton-density recovery and dictionary discretization error.
% Planned figures: Transient fingerprints and parameter ambiguity.
% Planned challenge problems: Construct indistinguishable fingerprints; compare MRF and conventional mapping at fixed time.
\chapter{Magnetic Resonance Fingerprinting}
\label{ch:12}

\section{Transient multiparametric encoding}
\label{sec:12:transient-multiparametric-encoding}
\subsection{Variable flip angles, TR and preparation}
\subsection{Bloch and EPG signal evolution}
\subsection{Sequence diversity versus parameter sensitivity}

\section{Dictionary inference}
\label{sec:12:dictionary-inference}
\subsection{Simulation grids and normalization}
\subsection{Matching metrics and proton-density scaling}
\subsection{Discretization, interpolation and nuisance parameters}

\section{Acquisition and reconstruction}
\label{sec:12:acquisition-and-reconstruction}
\subsection{Undersampling and temporal artifacts}
\subsection{Low-rank subspaces and compressed dictionaries}
\subsection{Model-based and iterative alternatives}

\section{Experiment design and uncertainty}
\label{sec:12:experiment-design-and-uncertainty}
\subsection{Identifiability and correlated fingerprints}
\subsection{Fisher information and sequence optimization}
\subsection{B0, B1-plus, motion and model mismatch}

\section{Validation and comparison}
\label{sec:12:validation-and-comparison}
\subsection{Repeatability and phantom validation}
\subsection{Fair scan-time and resolution comparisons}
\subsection{Clinical opportunities and unresolved limitations}

\section{Worked examples}
\label{sec:12:worked-examples}
% Planned calculations: Normalized matching, proton-density recovery and dictionary discretization error.
\subsection{Derivation and quantitative calculation}
\subsection{Assumptions, limiting cases and scanner interpretation}

\section{Challenge problems}
\label{sec:12:challenge-problems}
% Problem design: Construct indistinguishable fingerprints; compare MRF and conventional mapping at fixed time.
% Use challengeproblem and stable labels prob:12:<topic>.
% Complete solutions belong in Appendix E, section for this chapter.
% Use \begin{solution}{prob:12:<topic>} ... \end{solution} there.
\subsection{Derivation and quantitative design}
\subsection{Model criticism and integrated reasoning}

% Scope: Derive EPG from spatial Fourier coefficients and the declared rotation convention; track transverse and longitudinal states through RF, relaxation and gradients.
% Status: architecture only; no chapter prose or completed problems yet.
% Prerequisites: Chapters 3 through 8; Chapter 10 for diffusion attenuation.
% Planned worked examples: Stimulated echo, CPMG echo train and imperfect spoiling.
% Planned figures: Signed-order EPG state evolution and pathway diagrams.
% Planned challenge problems: Derive the RF transition matrix; verify EPG against isochromats.
\chapter{Extended Phase Graphs and Signal Pathways}
\label{ch:13}

\section{Configuration-state representation}
\label{sec:13:configuration-state-representation}
\subsection{Spatial harmonics and averaging}
\subsection{F-plus, F-minus and longitudinal Z states}
\subsection{Signed orders, conjugacy and observable echoes}

\section{EPG operators}
\label{sec:13:epg-operators}
\subsection{RF transition matrix from Bloch rotations}
\subsection{Relaxation, recovery and off-resonance}
\subsection{Gradient shifts and state truncation}

\section{Echo pathways}
\label{sec:13:echo-pathways}
\subsection{Hahn echoes and stimulated echoes}
\subsection{Crushers, spoiling and coherence selection}
\subsection{Phase cycling and pathway interference}

\section{Echo trains and steady states}
\label{sec:13:echo-trains-and-steady-states}
\subsection{CPMG and variable-angle FSE/TSE}
\subsection{Spoiled GRE and incomplete spoiling}
\subsection{SSFP and transient MRF applications}

\section{Extensions and numerical verification}
\label{sec:13:extensions-and-numerical-verification}
\subsection{Diffusion weighting along pathways}
\subsection{Slice profiles and noninteger dephasing}
\subsection{Convergence and comparison with Bloch ensembles}

\section{Worked examples}
\label{sec:13:worked-examples}
% Planned calculations: Stimulated echo, CPMG echo train and imperfect spoiling.
\subsection{Derivation and quantitative calculation}
\subsection{Assumptions, limiting cases and scanner interpretation}

\section{Challenge problems}
\label{sec:13:challenge-problems}
% Problem design: Derive the RF transition matrix; verify EPG against isochromats.
% Use challengeproblem and stable labels prob:13:<topic>.
% Complete solutions belong in Appendix E, section for this chapter.
% Use \begin{solution}{prob:13:<topic>} ... \end{solution} there.
\subsection{Derivation and quantitative design}
\subsection{Model criticism and integrated reasoning}


\part{Accelerated Imaging and Reconstruction}
% Scope: Derive coil encoding, SENSE, GRAPPA and noise amplification; separate acceleration geometry from reconstruction assumptions.
% Status: architecture only; no chapter prose or completed problems yet.
% Prerequisites: Chapters 1, 5 and 6; receive hardware in Chapter 17.
% Planned worked examples: Two-coil SENSE inversion and g-factor SNR penalty.
% Planned figures: Aliased voxels, sensitivity vectors and GRAPPA neighborhoods.
% Planned challenge problems: Find an encoding null space; compare calibration and noise penalties.
\chapter{Parallel Imaging}
\label{ch:14}

\section{Receive arrays as encoding systems}
\label{sec:14:receive-arrays-as-encoding-systems}
\subsection{Sensitivity fields and folded voxels}
\subsection{Noise covariance and prewhitening}
\subsection{Optimal coil combination and normalization}

\section{SENSE reconstruction}
\label{sec:14:sense-reconstruction}
\subsection{Local unfolding equations}
\subsection{Weighted least squares and conditioning}
\subsection{Noise propagation and g-factor}

\section{GRAPPA reconstruction}
\label{sec:14:grappa-reconstruction}
\subsection{Autocalibration and neighborhood kernels}
\subsection{Calibration fits and synthesized samples}
\subsection{Noise correlations and calibration failure}

\section{Simultaneous multislice reconstruction}
\label{sec:14:simultaneous-multislice-reconstruction}
\subsection{Slice separation and coil geometry}
\subsection{CAIPIRINHA shifts and controlled aliasing}
\subsection{Slice leakage and reconstruction trade-offs}

\section{Practical acceleration limits}
\label{sec:14:practical-acceleration-limits}
\subsection{Calibration overhead and net acceleration}
\subsection{Motion and sensitivity mismatch}
\subsection{Regularization, residual aliasing and SNR efficiency}

\section{Worked examples}
\label{sec:14:worked-examples}
% Planned calculations: Two-coil SENSE inversion and g-factor SNR penalty.
\subsection{Derivation and quantitative calculation}
\subsection{Assumptions, limiting cases and scanner interpretation}

\section{Challenge problems}
\label{sec:14:challenge-problems}
% Problem design: Find an encoding null space; compare calibration and noise penalties.
% Use challengeproblem and stable labels prob:14:<topic>.
% Complete solutions belong in Appendix E, section for this chapter.
% Use \begin{solution}{prob:14:<topic>} ... \end{solution} there.
\subsection{Derivation and quantitative design}
\subsection{Model criticism and integrated reasoning}

% Scope: Connect nonuniform sampling to gridding and operator-based reconstruction; explain acceleration as a joint sampling and prior-design problem.
% Status: architecture only; no chapter prose or completed problems yet.
% Prerequisites: Chapters 5, 6 and 14; solver details in Chapter 16.
% Planned worked examples: Radial sampling density, spiral off-resonance and adjoint checks.
% Planned figures: Trajectories, interpolation kernels and undersampling PSFs.
% Planned challenge problems: Design a trajectory under slew constraints; distinguish density compensation from inversion.
\chapter{Non-Cartesian and Accelerated Imaging}
\label{ch:15}

\section{Radial and spiral acquisition}
\label{sec:15:radial-and-spiral-acquisition}
\subsection{Angular sampling and golden-angle ordering}
\subsection{Spiral geometry and variable density}
\subsection{Stack-of-stars and three-dimensional trajectories}

\section{Nonuniform Fourier reconstruction}
\label{sec:15:nonuniform-fourier-reconstruction}
\subsection{Gridding kernels and oversampling}
\subsection{Deapodization and density compensation}
\subsection{NUFFT forward and adjoint operators}

\section{Trajectory and field errors}
\label{sec:15:trajectory-and-field-errors}
\subsection{Gradient delays and measured trajectories}
\subsection{Off-resonance blur and time segmentation}
\subsection{Motion robustness and self-navigation limits}

\section{Sampling for acceleration}
\label{sec:15:sampling-for-acceleration}
\subsection{Variable-density masks and incoherence}
\subsection{Partial Fourier with complex phase}
\subsection{Joint coil and sampling design}

\section{Dynamic imaging}
\label{sec:15:dynamic-imaging}
\subsection{Temporal interleaving and view sharing}
\subsection{Spatiotemporal undersampling and low-rank structure}
\subsection{Temporal footprint versus nominal frame rate}

\section{Worked examples}
\label{sec:15:worked-examples}
% Planned calculations: Radial sampling density, spiral off-resonance and adjoint checks.
\subsection{Derivation and quantitative calculation}
\subsection{Assumptions, limiting cases and scanner interpretation}

\section{Challenge problems}
\label{sec:15:challenge-problems}
% Problem design: Design a trajectory under slew constraints; distinguish density compensation from inversion.
% Use challengeproblem and stable labels prob:15:<topic>.
% Complete solutions belong in Appendix E, section for this chapter.
% Use \begin{solution}{prob:15:<topic>} ... \end{solution} there.
\subsection{Derivation and quantitative design}
\subsection{Model criticism and integrated reasoning}

% Scope: Build an operator-based account of reconstruction from least squares to regularization and model-based inference; restrict learned methods to physics, stability and validation.
% Status: architecture only; no chapter prose or completed problems yet.
% Prerequisites: Chapters 1, 6 and 14 through 15.
% Planned worked examples: Tikhonov solution, adjoint consistency and iterative convergence.
% Planned figures: Singular spectra, bias-variance curves and reconstruction residuals.
% Planned challenge problems: Prove nonuniqueness; compare reconstructions with matched data fidelity.
\chapter{Reconstruction and Inverse Problems}
\label{ch:16}

\section{The discrete forward model}
\label{sec:16:the-discrete-forward-model}
\subsection{Encoding operators and discretization}
\subsection{Adjoint versus inverse and inner-product tests}
\subsection{Noise whitening and likelihood}

\section{Linear reconstruction}
\label{sec:16:linear-reconstruction}
\subsection{Least squares, pseudoinverses and rank}
\subsection{Tikhonov regularization and singular-value filtering}
\subsection{Conditioning, resolution matrices and uncertainty}

\section{Sparsity and compressed sensing}
\label{sec:16:sparsity-and-compressed-sensing}
\subsection{Transform sparsity and incoherent measurements}
\subsection{L1 and total-variation objectives}
\subsection{Sampling assumptions and failure cases}

\section{Iterative algorithms}
\label{sec:16:iterative-algorithms}
\subsection{Conjugate gradients and normal equations}
\subsection{Proximal gradients and variable splitting}
\subsection{Stopping rules, preconditioning and data consistency}

\section{Low-rank and model-based reconstruction}
\label{sec:16:low-rank-and-model-based-reconstruction}
\subsection{Dynamic subspaces and low-rank penalties}
\subsection{Joint image and parameter estimation}
\subsection{Nonconvexity, initialization and identifiability}

\section{Learned reconstruction}
\label{sec:16:learned-reconstruction}
\subsection{Learned priors and unrolled optimization}
\subsection{Distribution shift, hallucination and uncertainty}
\subsection{Validation with residuals and task-based metrics}

\section{Worked examples}
\label{sec:16:worked-examples}
% Planned calculations: Tikhonov solution, adjoint consistency and iterative convergence.
\subsection{Derivation and quantitative calculation}
\subsection{Assumptions, limiting cases and scanner interpretation}

\section{Challenge problems}
\label{sec:16:challenge-problems}
% Problem design: Prove nonuniqueness; compare reconstructions with matched data fidelity.
% Use challengeproblem and stable labels prob:16:<topic>.
% Complete solutions belong in Appendix E, section for this chapter.
% Use \begin{solution}{prob:16:<topic>} ... \end{solution} there.
\subsection{Derivation and quantitative design}
\subsection{Model criticism and integrated reasoning}


\part{Hardware, Sequence Design and Specialized Measurements}
% Scope: Explain the magnet, gradient and RF subsystems as physical implementations of earlier models; connect imperfections to encoding and safety limits.
% Status: architecture only; no chapter prose or completed problems yet.
% Prerequisites: Chapters 1, 4 through 6 and 14.
% Planned worked examples: Gradient slew, receive-noise budget and digitization timing.
% Planned figures: Magnet cross-section, Maxwell/Golay coils and receiver chain.
% Planned challenge problems: Trace a hardware fault through k-space; allocate a receiver noise budget.
\chapter{MRI Hardware}
\label{ch:17}

\section{Main magnet and field control}
\label{sec:17:main-magnet-and-field-control}
\subsection{Superconductivity, persistent currents and cryogenics}
\subsection{Homogeneity, drift and passive shimming}
\subsection{Active shims and field monitoring}

\section{Gradient subsystem}
\label{sec:17:gradient-subsystem}
\subsection{Maxwell pairs, Golay coils and linearity}
\subsection{Amplifiers, inductance, amplitude and slew}
\subsection{Eddy currents, mechanical vibration and PNS links}

\section{RF transmit subsystem}
\label{sec:17:rf-transmit-subsystem}
\subsection{Birdcage coils, quadrature and loading}
\subsection{B1-plus fields and transmit calibration}
\subsection{Parallel transmission and coupling}

\section{RF receive subsystem}
\label{sec:17:rf-receive-subsystem}
\subsection{Surface coils and phased arrays}
\subsection{Coil sensitivity, sample noise and preamplifiers}
\subsection{Decoupling, matching and transmit-receive protection}

\section{Receiver and control chain}
\label{sec:17:receiver-and-control-chain}
\subsection{Mixing, filtering and complex demodulation}
\subsection{ADC bandwidth, dynamic range and quantization}
\subsection{Timing, synchronization and gradient monitoring}

\section{System-level limitations}
\label{sec:17:system-level-limitations}
\subsection{Field strength and noise regimes}
\subsection{Calibration drift and quality assurance}
\subsection{Hardware constraints in sequence design}

\section{Worked examples}
\label{sec:17:worked-examples}
% Planned calculations: Gradient slew, receive-noise budget and digitization timing.
\subsection{Derivation and quantitative calculation}
\subsection{Assumptions, limiting cases and scanner interpretation}

\section{Challenge problems}
\label{sec:17:challenge-problems}
% Problem design: Trace a hardware fault through k-space; allocate a receiver noise budget.
% Use challengeproblem and stable labels prob:17:<topic>.
% Complete solutions belong in Appendix E, section for this chapter.
% Use \begin{solution}{prob:17:<topic>} ... \end{solution} there.
\subsection{Derivation and quantitative design}
\subsection{Model criticism and integrated reasoning}

% Scope: Develop selective excitation and constrained sequence synthesis using the established RF conventions; distinguish mathematical designs from scanner-feasible waveforms.
% Status: architecture only; no chapter prose or completed problems yet.
% Prerequisites: Chapters 4, 5, 8 and 17.
% Planned worked examples: Time-bandwidth product, slice thickness and minimum gradient timing.
% Planned figures: Sinc profiles, excitation k-space and gradient ramps.
% Planned challenge problems: Design a slice-selective pulse; compare SAR and duration under hardware limits.
\chapter{RF Pulses and Sequence Design}
\label{ch:18}

\section{Selective excitation theory}
\label{sec:18:selective-excitation-theory}
\subsection{Small-tip-angle integral and excitation k-space}
\subsection{Sinc pulses, apodization and slice profiles}
\subsection{RF bandwidth and time-bandwidth product}

\section{Large-tip and robust pulses}
\label{sec:18:large-tip-and-robust-pulses}
\subsection{Composite rotations and error compensation}
\subsection{Adiabatic following and inversion}
\subsection{SLR design and excitation-refocusing differences}

\section{Multidimensional and multiband excitation}
\label{sec:18:multidimensional-and-multiband-excitation}
\subsection{Spatially selective trajectories}
\subsection{Multiband phases and peak RF power}
\subsection{Parallel transmit design and local heating constraints}

\section{Joint RF and gradient design}
\label{sec:18:joint-rf-and-gradient-design}
\subsection{VERSE and waveform reparameterization}
\subsection{Off-resonance and relaxation limitations}
\subsection{Gradient amplitude, slew and raster constraints}

\section{Sequence timing and gradient moments}
\label{sec:18:sequence-timing-and-gradient-moments}
\subsection{Trapezoid design and minimum duration}
\subsection{Rephasing, crushers and flow compensation}
\subsection{TE, TR and echo-spacing feasibility}

\section{Design verification}
\label{sec:18:design-verification}
\subsection{Bloch simulation across position and detuning}
\subsection{Slice-profile and pathway validation}
\subsection{SAR implications and scanner implementation checks}

\section{Worked examples}
\label{sec:18:worked-examples}
% Planned calculations: Time-bandwidth product, slice thickness and minimum gradient timing.
\subsection{Derivation and quantitative calculation}
\subsection{Assumptions, limiting cases and scanner interpretation}

\section{Challenge problems}
\label{sec:18:challenge-problems}
% Problem design: Design a slice-selective pulse; compare SAR and duration under hardware limits.
% Use challengeproblem and stable labels prob:18:<topic>.
% Complete solutions belong in Appendix E, section for this chapter.
% Use \begin{solution}{prob:18:<topic>} ... \end{solution} there.
\subsection{Derivation and quantitative design}
\subsection{Model criticism and integrated reasoning}

% Scope: Derive motion-induced phase and inflow contrast; develop phase contrast, TOF and ASL with explicit timing, calibration and physiological assumptions.
% Status: architecture only; no chapter prose or completed problems yet.
% Prerequisites: Chapters 5, 7 through 9 and 17.
% Planned worked examples: Gradient moments, velocity encoding and ASL timing.
% Planned figures: Bipolar gradients, inflow saturation and label-control timing.
% Planned challenge problems: Design a VENC experiment; distinguish transit-time bias from low perfusion.
\chapter{Flow, Perfusion and Angiography}
\label{ch:19}

\section{Motion in gradient fields}
\label{sec:19:motion-in-gradient-fields}
\subsection{Position expansion and gradient moments}
\subsection{Constant velocity, acceleration and intravoxel dispersion}
\subsection{Flow compensation and residual sensitivity}

\section{Phase-contrast imaging}
\label{sec:19:phase-contrast-imaging}
\subsection{Bipolar encoding and velocity phase}
\subsection{VENC, phase wrapping and noise}
\subsection{Background phase and multi-directional flow}

\section{Inflow and angiographic contrast}
\label{sec:19:inflow-and-angiographic-contrast}
\subsection{Time-of-flight saturation and fresh spins}
\subsection{Slice/slab geometry and venous suppression}
\subsection{Contrast-enhanced angiography and timing}

\section{Arterial spin labeling}
\label{sec:19:arterial-spin-labeling}
\subsection{Pulsed, continuous and pseudocontinuous labeling}
\subsection{Label-control subtraction and kinetic models}
\subsection{Transit time, labeling efficiency and quantification}

\section{Applications and limitations}
\label{sec:19:applications-and-limitations}
\subsection{Vascular stenosis and complex flow}
\subsection{Cardiac gating and four-dimensional flow}
\subsection{Perfusion versus macrovascular contamination}

\section{Worked examples}
\label{sec:19:worked-examples}
% Planned calculations: Gradient moments, velocity encoding and ASL timing.
\subsection{Derivation and quantitative calculation}
\subsection{Assumptions, limiting cases and scanner interpretation}

\section{Challenge problems}
\label{sec:19:challenge-problems}
% Problem design: Design a VENC experiment; distinguish transit-time bias from low perfusion.
% Use challengeproblem and stable labels prob:19:<topic>.
% Complete solutions belong in Appendix E, section for this chapter.
% Use \begin{solution}{prob:19:<topic>} ... \end{solution} there.
\subsection{Derivation and quantitative design}
\subsection{Model criticism and integrated reasoning}

% Scope: Connect spin interactions and localization to measured spectra, including acquisition constraints and fitting ambiguity.
% Status: architecture only; no chapter prose or completed problems yet.
% Prerequisites: Chapters 2, 4, 6 and 9.
% Planned worked examples: Chemical shift in ppm, spectral resolution and linewidth.
% Planned figures: Coupled levels, spectra, PRESS and STEAM timing.
% Planned challenge problems: Predict echo-time-dependent multiplets; diagnose localization and baseline errors.
\chapter{Magnetic Resonance Spectroscopy}
\label{ch:20}

\section{Chemical shifts and spin interactions}
\label{sec:20:chemical-shifts-and-spin-interactions}
\subsection{Shielding and ppm references}
\subsection{J-coupling and coupled-spin Hamiltonians}
\subsection{Weak coupling, multiplets and evolution}

\section{From FID to spectrum}
\label{sec:20:from-fid-to-spectrum}
\subsection{Spectral Fourier transform and frequency sign}
\subsection{Linewidth, T2-star and apodization}
\subsection{Dwell time, spectral width and frequency resolution}

\section{Localization}
\label{sec:20:localization}
\subsection{Slice-selective localization and chemical-shift displacement}
\subsection{PRESS spin echoes}
\subsection{STEAM stimulated echoes and mixing time}

\section{Acquisition practice}
\label{sec:20:acquisition-practice}
\subsection{Shimming and field stability}
\subsection{Water and lipid suppression}
\subsection{Single-voxel and spectroscopic imaging}

\section{Spectral modeling and interpretation}
\label{sec:20:spectral-modeling-and-interpretation}
\subsection{Basis spectra, phase and baseline}
\subsection{Metabolites, concentration and reference scaling}
\subsection{Fitting uncertainty, artifacts and biological specificity}

\section{Worked examples}
\label{sec:20:worked-examples}
% Planned calculations: Chemical shift in ppm, spectral resolution and linewidth.
\subsection{Derivation and quantitative calculation}
\subsection{Assumptions, limiting cases and scanner interpretation}

\section{Challenge problems}
\label{sec:20:challenge-problems}
% Problem design: Predict echo-time-dependent multiplets; diagnose localization and baseline errors.
% Use challengeproblem and stable labels prob:20:<topic>.
% Complete solutions belong in Appendix E, section for this chapter.
% Use \begin{solution}{prob:20:<topic>} ... \end{solution} there.
\subsection{Derivation and quantitative design}
\subsection{Model criticism and integrated reasoning}

% Scope: Develop BOLD signal formation and acquisition trade-offs before a limited statistical treatment; distinguish vascular inference from direct neuronal measurement.
% Status: architecture only; no chapter prose or completed problems yet.
% Prerequisites: Chapters 8, 9 and 16.
% Planned worked examples: TE dependence, temporal sampling and physiological-noise limits.
% Planned figures: Hemodynamic response and temporal sampling.
% Planned challenge problems: Optimize an EPI protocol under dropout constraints; diagnose a confounded activation.
\chapter{Functional MRI}
\label{ch:21}

\section{BOLD contrast physics}
\label{sec:21:bold-contrast-physics}
\subsection{Deoxyhemoglobin and susceptibility}
\subsection{Intravascular and extravascular contributions}
\subsection{Gradient-echo and spin-echo sensitivity}

\section{Neurovascular coupling}
\label{sec:21:neurovascular-coupling}
\subsection{Blood flow, volume and oxygen metabolism}
\subsection{Hemodynamic response and temporal variability}
\subsection{Vascular specificity and limitations}

\section{Acquisition design}
\label{sec:21:acquisition-design}
\subsection{EPI, TE and echo spacing}
\subsection{Spatial resolution, temporal resolution and SMS}
\subsection{Distortion, dropout and physiological noise}

\section{Time-series inference}
\label{sec:21:time-series-inference}
\subsection{Drift, motion and nuisance regressors}
\subsection{General linear model and correlated noise}
\subsection{Multiple comparisons and statistical interpretation}

\section{Limits and validation}
\label{sec:21:limits-and-validation}
\subsection{Confounding and causal claims}
\subsection{Task and resting-state measurements}
\subsection{Reproducibility and acquisition-dependent bias}

\section{Worked examples}
\label{sec:21:worked-examples}
% Planned calculations: TE dependence, temporal sampling and physiological-noise limits.
\subsection{Derivation and quantitative calculation}
\subsection{Assumptions, limiting cases and scanner interpretation}

\section{Challenge problems}
\label{sec:21:challenge-problems}
% Problem design: Optimize an EPI protocol under dropout constraints; diagnose a confounded activation.
% Use challengeproblem and stable labels prob:21:<topic>.
% Complete solutions belong in Appendix E, section for this chapter.
% Use \begin{solution}{prob:21:<topic>} ... \end{solution} there.
\subsection{Derivation and quantitative design}
\subsection{Model criticism and integrated reasoning}


\part{Artifacts, Protocols and Integrated Reasoning}
% Scope: Use a uniform causal account for each artifact: physics, mathematical model, appearance, parameter dependence, diagnosis, mitigation and cost.
% Status: architecture only; no chapter prose or completed problems yet.
% Prerequisites: Chapters 5 through 9 and 14 through 18.
% Planned worked examples: Chemical-shift displacement, EPI distortion and ghost offsets.
% Planned figures: Artifact-producing k-space errors and corrected/uncorrected PSFs.
% Planned challenge problems: Infer a fault from acquisition settings; choose between imperfect correction strategies.
\chapter{Artifacts and Corrections}
\label{ch:22}

\section{Motion and flow artifacts}
\label{sec:22:motion-and-flow-artifacts}
\subsection{Inconsistent phase encoding and ghosting}
\subsection{Rigid motion, deformation and intra-readout motion}
\subsection{Navigators, gating and prospective correction}

\section{EPI-specific errors}
\label{sec:22:epi-specific-errors}
\subsection{Odd-even mismatch and Nyquist ghosts}
\subsection{Off-resonance geometric distortion and pile-up}
\subsection{Signal dropout and correction limits}

\section{Chemical shift and field inhomogeneity}
\label{sec:22:chemical-shift-and-field-inhomogeneity}
\subsection{Frequency displacement and fat-water cancellation}
\subsection{Susceptibility, B0 drift and shimming}
\subsection{B1 inhomogeneity, dielectric shading and fat-suppression failure}

\section{Sampling and finite resolution}
\label{sec:22:sampling-and-finite-resolution}
\subsection{Aliasing and wraparound}
\subsection{Gibbs ringing and truncation}
\subsection{Partial volume and interpolation artifacts}

\section{Hardware-related errors}
\label{sec:22:hardware-related-errors}
\subsection{RF interference and zipper artifacts}
\subsection{Eddy currents and trajectory delays}
\subsection{Gradient nonlinearity and concomitant fields}

\section{Reconstruction-related errors}
\label{sec:22:reconstruction-related-errors}
\subsection{Parallel-imaging residual aliasing and noise}
\subsection{Non-Cartesian streaking and blur}
\subsection{Regularization, learned priors and synthetic structure}

\section{Diagnosis and correction trade-offs}
\label{sec:22:diagnosis-and-correction-trade-offs}
\subsection{Parameter perturbations as controlled experiments}
\subsection{Correction-induced noise and resolution loss}
\subsection{Validation against raw data and independent acquisitions}

\section{Worked examples}
\label{sec:22:worked-examples}
% Planned calculations: Chemical-shift displacement, EPI distortion and ghost offsets.
\subsection{Derivation and quantitative calculation}
\subsection{Assumptions, limiting cases and scanner interpretation}

\section{Challenge problems}
\label{sec:22:challenge-problems}
% Problem design: Infer a fault from acquisition settings; choose between imperfect correction strategies.
% Use challengeproblem and stable labels prob:22:<topic>.
% Complete solutions belong in Appendix E, section for this chapter.
% Use \begin{solution}{prob:22:<topic>} ... \end{solution} there.
\subsection{Derivation and quantitative design}
\subsection{Model criticism and integrated reasoning}

% Scope: Synthesize sequence and parameter choices by clinical objective, with anatomy used for examples rather than the book's organizing principle.
% Status: architecture only; no chapter prose or completed problems yet.
% Prerequisites: Chapters 7 through 11 and 19 through 22.
% Planned worked examples: Brain FLAIR, cardiac cine/LGE and liver fat-water protocol comparisons.
% Planned figures: Protocol decision maps tied to physical contrast.
% Planned challenge problems: Design a fixed-time protocol; justify every trade-off and identify failure modes.
\chapter{Clinical Protocol Physics}
\label{ch:23}

\section{From diagnostic task to acquisition}
\label{sec:23:from-diagnostic-task-to-acquisition}
\subsection{Target contrast and spatial scale}
\subsection{Motion, timing and patient constraints}
\subsection{Redundancy, uncertainty and protocol robustness}

\section{Brain protocols}
\label{sec:23:brain-protocols}
\subsection{Structural MPRAGE and FLAIR suppression}
\subsection{Acute ischemia and diffusion interpretation}
\subsection{Susceptibility, hemorrhage and T2-star}

\section{Cardiac protocols}
\label{sec:23:cardiac-protocols}
\subsection{bSSFP cine and off-resonance management}
\subsection{Gating, motion and temporal resolution}
\subsection{Late gadolinium enhancement and inversion timing}

\section{Liver and abdominal protocols}
\label{sec:23:liver-and-abdominal-protocols}
\subsection{Breath holding and respiratory motion}
\subsection{Dixon, fat fraction and iron-sensitive decay}
\subsection{Diffusion and dynamic contrast timing}

\section{Musculoskeletal and oncologic protocols}
\label{sec:23:musculoskeletal-and-oncologic-protocols}
\subsection{Fat suppression and fluid-sensitive imaging}
\subsection{Spatial resolution and partial-volume control}
\subsection{Multiparametric tumor assessment and specificity}

\section{Vascular protocol synthesis}
\label{sec:23:vascular-protocol-synthesis}
\subsection{TOF, phase contrast and contrast enhancement}
\subsection{Velocity range and flow artifacts}
\subsection{Choosing coverage, resolution and scan time}

\section{Worked examples}
\label{sec:23:worked-examples}
% Planned calculations: Brain FLAIR, cardiac cine/LGE and liver fat-water protocol comparisons.
\subsection{Derivation and quantitative calculation}
\subsection{Assumptions, limiting cases and scanner interpretation}

\section{Challenge problems}
\label{sec:23:challenge-problems}
% Problem design: Design a fixed-time protocol; justify every trade-off and identify failure modes.
% Use challengeproblem and stable labels prob:23:<topic>.
% Complete solutions belong in Appendix E, section for this chapter.
% Use \begin{solution}{prob:23:<topic>} ... \end{solution} there.
\subsection{Derivation and quantitative design}
\subsection{Model criticism and integrated reasoning}

% Scope: Quantify coupled parameter trade-offs with acquisition and reconstruction assumptions stated; optimize task performance rather than isolated image appearance.
% Status: architecture only; no chapter prose or completed problems yet.
% Prerequisites: Chapters 6 through 16 and 22 through 23.
% Planned worked examples: SNR scaling, bandwidth displacement, scan time and acceleration penalty.
% Planned figures: Pareto curves, SNR efficiency and contrast surfaces.
% Planned challenge problems: Optimize a fixed-time protocol; explain why a familiar SNR scaling fails.
\chapter{Image Quality and Optimization}
\label{ch:24}

\section{Image-quality measures}
\label{sec:24:image-quality-measures}
\subsection{SNR, CNR and noise definitions}
\subsection{PSF, effective resolution and detectability}
\subsection{Repeatability and task-based assessment}

\section{Sampling and averaging trade-offs}
\label{sec:24:sampling-and-averaging-trade-offs}
\subsection{FOV, matrix, voxel volume and slice thickness}
\subsection{Number of averages and scan-time scaling}
\subsection{Two-dimensional versus three-dimensional efficiency}

\section{Timing and contrast trade-offs}
\label{sec:24:timing-and-contrast-trade-offs}
\subsection{TR, TE, TI and flip angle}
\subsection{Echo train length, turbo factor and effective TE}
\subsection{Contrast, blurring and motion sensitivity}

\section{Bandwidth and gradient trade-offs}
\label{sec:24:bandwidth-and-gradient-trade-offs}
\subsection{Receiver bandwidth, dwell time and noise}
\subsection{RF bandwidth and slice selection}
\subsection{Echo spacing, readout duration, distortion and slew}

\section{Acceleration and incomplete sampling}
\label{sec:24:acceleration-and-incomplete-sampling}
\subsection{Parallel-imaging SNR and calibration overhead}
\subsection{Partial Fourier and phase uncertainty}
\subsection{SMS, compressed sensing and nonlinear image quality}

\section{Constrained optimization}
\label{sec:24:constrained-optimization}
\subsection{Fixed-time versus fixed-resolution comparisons}
\subsection{SAR, motion and hardware constraints}
\subsection{Pareto design and sensitivity to model assumptions}

\section{Worked examples}
\label{sec:24:worked-examples}
% Planned calculations: SNR scaling, bandwidth displacement, scan time and acceleration penalty.
\subsection{Derivation and quantitative calculation}
\subsection{Assumptions, limiting cases and scanner interpretation}

\section{Challenge problems}
\label{sec:24:challenge-problems}
% Problem design: Optimize a fixed-time protocol; explain why a familiar SNR scaling fails.
% Use challengeproblem and stable labels prob:24:<topic>.
% Complete solutions belong in Appendix E, section for this chapter.
% Use \begin{solution}{prob:24:<topic>} ... \end{solution} there.
\subsection{Derivation and quantitative design}
\subsection{Model criticism and integrated reasoning}

% Scope: Integrate mature modern methods through their physical opportunities and limitations; point back to full derivations and avoid a catalogue of recent papers.
% Status: architecture only; no chapter prose or completed problems yet.
% Prerequisites: Chapters 11 through 18 and 24.
% Planned worked examples: Field-strength scaling and comparison of quantitative acquisition strategies.
% Planned figures: Method dependency map and low-field trade-offs.
% Planned challenge problems: Compare low-field and high-field task performance; audit a learned reconstruction claim.
\chapter{Modern MRI Methods}
\label{ch:25}

\section{Low-field and portable MRI}
\label{sec:25:low-field-and-portable-mri}
\subsection{Polarization, relaxation and noise regimes}
\subsection{Magnet design, bandwidth and encoding constraints}
\subsection{Resolution, accessibility and reconstruction trade-offs}

\section{High-field and parallel transmission}
\label{sec:25:high-field-and-parallel-transmission}
\subsection{Wavelength effects and B1-plus control}
\subsection{Joint pulse and field design}
\subsection{SAR, susceptibility and calibration constraints}

\section{Short-lived and alternative signals}
\label{sec:25:short-lived-and-alternative-signals}
\subsection{Ultrashort- and zero-echo-time acquisition}
\subsection{Nonproton and hyperpolarized imaging concepts}
\subsection{Hardware and quantification limitations}

\section{Integrated quantitative imaging}
\label{sec:25:integrated-quantitative-imaging}
\subsection{Synthetic MRI and multiparametric acquisitions}
\subsection{MRF and model-based reconstruction in context}
\subsection{Calibration, harmonization and uncertainty}

\section{Modern reconstruction in practice}
\label{sec:25:modern-reconstruction-in-practice}
\subsection{Compressed sensing, low-rank and learned methods}
\subsection{Prospective validation and distribution shift}
\subsection{Evidence maturity and reproducible comparisons}

\section{Worked examples}
\label{sec:25:worked-examples}
% Planned calculations: Field-strength scaling and comparison of quantitative acquisition strategies.
\subsection{Derivation and quantitative calculation}
\subsection{Assumptions, limiting cases and scanner interpretation}

\section{Challenge problems}
\label{sec:25:challenge-problems}
% Problem design: Compare low-field and high-field task performance; audit a learned reconstruction claim.
% Use challengeproblem and stable labels prob:25:<topic>.
% Complete solutions belong in Appendix E, section for this chapter.
% Use \begin{solution}{prob:25:<topic>} ... \end{solution} there.
\subsection{Derivation and quantitative design}
\subsection{Model criticism and integrated reasoning}

% Scope: Explain static, RF, gradient, acoustic and cryogenic hazards from their mechanisms; keep operational limits tied to current authoritative standards when drafted.
% Status: architecture only; no chapter prose or completed problems yet.
% Prerequisites: Chapters 1, 17, 18 and 24.
% Planned worked examples: RF duty-cycle scaling and gradient slew comparisons without patient-specific clearance.
% Planned figures: Energy deposition and induced electric-field pathways.
% Planned challenge problems: Compare equal-flip pulses with unequal heating; distinguish SAR from local temperature.
\chapter{Safety and System Limits}
\label{ch:26}

\section{Static-field interactions}
\label{sec:26:static-field-interactions}
\subsection{Forces, torques and projectile mechanisms}
\subsection{Spatial gradients and active shielding}
\subsection{Implant interactions and screening principles}

\section{RF energy deposition}
\label{sec:26:rf-energy-deposition}
\subsection{Electric fields, dissipated power and SAR}
\subsection{Global versus local heating and duty cycle}
\subsection{Conductive loops, implants and model uncertainty}

\section{Gradient and acoustic effects}
\label{sec:26:gradient-and-acoustic-effects}
\subsection{Induced electric fields and PNS}
\subsection{Waveform dependence beyond peak slew}
\subsection{Mechanical forces, acoustic noise and protection}

\section{Cryogenics and system failures}
\label{sec:26:cryogenics-and-system-failures}
\subsection{Stored energy and quench mechanisms}
\subsection{Cryogen release and oxygen displacement}
\subsection{Fault response and system interlocks}

\section{Limits and responsible interpretation}
\label{sec:26:limits-and-responsible-interpretation}
\subsection{Scanner operating modes and standards}
\subsection{MR Safe, MR Conditional and MR Unsafe labeling}
\subsection{Physics estimates versus device-specific conditions}

\section{Worked examples}
\label{sec:26:worked-examples}
% Planned calculations: RF duty-cycle scaling and gradient slew comparisons without patient-specific clearance.
\subsection{Derivation and quantitative calculation}
\subsection{Assumptions, limiting cases and scanner interpretation}

\section{Challenge problems}
\label{sec:26:challenge-problems}
% Problem design: Compare equal-flip pulses with unequal heating; distinguish SAR from local temperature.
% Use challengeproblem and stable labels prob:26:<topic>.
% Complete solutions belong in Appendix E, section for this chapter.
% Use \begin{solution}{prob:26:<topic>} ... \end{solution} there.
\subsection{Derivation and quantitative design}
\subsection{Model criticism and integrated reasoning}

% Scope: Reserve substantial synthesis problems and worked cases integrating spin physics, design, reconstruction and clinical objectives; solutions to challenge problems remain in Appendix E.
% Status: architecture only; no chapter prose or completed problems yet.
% Prerequisites: All preceding chapters as indicated per problem.
% Planned worked examples: End-to-end scanner design and quantitative protocol audits.
% Planned figures: Complete timing diagrams and simulated failure cases.
% Planned challenge problems: Multi-stage derivations, constrained design, inverse problems and research interview defenses.
\chapter{Integrated Worked Problems}
\label{ch:27}

\section{End-to-end acquisition design}
\label{sec:27:end-to-end-acquisition-design}
\subsection{From target resolution to gradient waveform}
\subsection{Sequence timing, contrast and scan time}
\subsection{Hardware and safety feasibility}

\section{Signal and pathway prediction}
\label{sec:27:signal-and-pathway-prediction}
\subsection{Imperfect pulses and stimulated echoes}
\subsection{Transient preparation and steady-state limits}
\subsection{Comparing Bloch, EPG and measured signals}

\section{Reconstruction and identifiability}
\label{sec:27:reconstruction-and-identifiability}
\subsection{Joint coil, field and image estimation}
\subsection{Undersampling ambiguity and prior dependence}
\subsection{Noise amplification and uncertainty bounds}

\section{Quantitative experimental design}
\label{sec:27:quantitative-experimental-design}
\subsection{Competing relaxation and diffusion models}
\subsection{Sampling schedules and nuisance parameters}
\subsection{Bias-variance trade-offs under fixed time}

\section{Artifact diagnosis and protocol repair}
\label{sec:27:artifact-diagnosis-and-protocol-repair}
\subsection{Coupled field, motion and sampling errors}
\subsection{Controlled tests and competing hypotheses}
\subsection{Correction limits and remaining uncertainty}

\section{Research and interview synthesis}
\label{sec:27:research-and-interview-synthesis}
\subsection{Derive an unfamiliar sequence from its diagram}
\subsection{Critique an apparently convincing result}
\subsection{Design a falsifiable MRI experiment}

\section{Worked examples}
\label{sec:27:worked-examples}
% Planned calculations: End-to-end scanner design and quantitative protocol audits.
\subsection{Derivation and quantitative calculation}
\subsection{Assumptions, limiting cases and scanner interpretation}

\section{Challenge problems}
\label{sec:27:challenge-problems}
% Problem design: Multi-stage derivations, constrained design, inverse problems and research interview defenses.
% Use challengeproblem and stable labels prob:27:<topic>.
% Complete solutions belong in Appendix E, section for this chapter.
% Use \begin{solution}{prob:27:<topic>} ... \end{solution} there.
\subsection{Derivation and quantitative design}
\subsection{Model criticism and integrated reasoning}


\printbibliography[heading=bibintoc,title={References}]
\appendix
% Scope: Collect reusable identities with assumptions and links to their derivations; do not replace Chapter 1's development.
% Status: architecture only; no chapter prose or completed problems yet.
\chapter{Mathematical Reference}
\label{app:a}

\section{Complex analysis and vector identities}
\label{sec:a:complex-analysis-and-vector-identities}
\subsection{Phasors and complex differentiation}
\subsection{Vector products and rotation identities}

\section{Linear operators and matrix calculus}
\label{sec:a:linear-operators-and-matrix-calculus}
\subsection{Spectral decompositions and matrix exponentials}
\subsection{Gradients, Jacobians and Hessians}

\section{Probability and estimation identities}
\label{sec:a:probability-and-estimation-identities}
\subsection{Gaussian likelihoods and covariance propagation}
\subsection{Fisher information and constrained estimation}

\section{Differential equations and propagators}
\label{sec:a:differential-equations-and-propagators}
\subsection{Linear systems and affine evolution}
\subsection{Numerical accuracy and limiting cases}

% Scope: Provide transform pairs and discrete conventions consistent with the front matter, including normalization and units.
% Status: architecture only; no chapter prose or completed problems yet.
\chapter{Fourier Reference}
\label{app:b}

\section{Continuous transform pairs}
\label{sec:b:continuous-transform-pairs}
\subsection{Gaussian, rectangle and sinc}
\subsection{Shifts, modulation and convolution}

\section{Discrete transforms and sampling}
\label{sec:b:discrete-transforms-and-sampling}
\subsection{DFT normalization and index ordering}
\subsection{Sampling combs and aliasing}

\section{Imaging and spectral conventions}
\label{sec:b:imaging-and-spectral-conventions}
\subsection{Spatial forward and inverse transforms}
\subsection{Temporal spectra and signed frequency axes}

\section{Resolution and windowing}
\label{sec:b:resolution-and-windowing}
\subsection{Point-spread functions and apodization}
\subsection{Zero filling and interpolation}

% Scope: Collect validated signal models only after their derivations exist, with assumptions, units and failure modes next to each equation.
% Status: architecture only; no chapter prose or completed problems yet.
\chapter{Common Signal Equations}
\label{app:c}

\section{Equilibrium and elementary evolution}
\label{sec:c:equilibrium-and-elementary-evolution}
\subsection{Thermal magnetization and Bloch free evolution}
\subsection{FID, spin echo and gradient echo}

\section{Steady-state and prepared acquisitions}
\label{sec:c:steady-state-and-prepared-acquisitions}
\subsection{Spoiled GRE and bSSFP}
\subsection{Inversion recovery and prepared readouts}

\section{Quantitative signal families}
\label{sec:c:quantitative-signal-families}
\subsection{Relaxometry, diffusion and exchange}
\subsection{Flow, ASL and chemical-species models}

\section{Noise and estimation}
\label{sec:c:noise-and-estimation}
\subsection{Complex likelihoods and magnitude statistics}
\subsection{SNR scaling and uncertainty relations}

% Scope: Separate physical constants from field-, tissue-, vendor- and protocol-dependent representative values; record sources and measurement conditions.
% Status: architecture only; no chapter prose or completed problems yet.
\chapter{Constants, Units and Typical Values}
\label{app:d}

\section{Constants and unit conversions}
\label{sec:d:constants-and-unit-conversions}
\subsection{Gyromagnetic ratios and fundamental constants}
\subsection{SI conversions and cycles versus radians}

\section{Tissue and field-dependent quantities}
\label{sec:d:tissue-and-field-dependent-quantities}
\subsection{Relaxation and diffusion ranges}
\subsection{Susceptibility and chemical shifts}

\section{System and acquisition quantities}
\label{sec:d:system-and-acquisition-quantities}
\subsection{Gradients, slew, bandwidth and timing}
\subsection{Nominal values versus measured performance}

% Scope: Provide complete derivations and numerical checks keyed to stable problem labels; keep all challenge solutions separate from questions.
% Status: architecture only; no chapter prose or completed problems yet.
\chapter{Solutions to Challenge Problems}
\label{app:e}

\section{Mathematical and Physical Foundations}
\label{sec:e:mathematical-and-physical-foundations}
\subsection{Derivations and numerical checks}
\subsection{Interpretation, limiting cases and alternative approaches}

\section{Quantum Spin Physics}
\label{sec:e:quantum-spin-physics}
\subsection{Derivations and numerical checks}
\subsection{Interpretation, limiting cases and alternative approaches}

\section{Classical Magnetization and the Bloch Equations}
\label{sec:e:classical-magnetization-and-the-bloch-equations}
\subsection{Derivations and numerical checks}
\subsection{Interpretation, limiting cases and alternative approaches}

\section{RF Excitation and the Rotating Frame}
\label{sec:e:rf-excitation-and-the-rotating-frame}
\subsection{Derivations and numerical checks}
\subsection{Interpretation, limiting cases and alternative approaches}

\section{Spatial Encoding and k-Space}
\label{sec:e:spatial-encoding-and-k-space}
\subsection{Derivations and numerical checks}
\subsection{Interpretation, limiting cases and alternative approaches}

\section{Signal Formation and Fourier Imaging}
\label{sec:e:signal-formation-and-fourier-imaging}
\subsection{Derivations and numerical checks}
\subsection{Interpretation, limiting cases and alternative approaches}

\section{Basic Pulse Sequences}
\label{sec:e:basic-pulse-sequences}
\subsection{Derivations and numerical checks}
\subsection{Interpretation, limiting cases and alternative approaches}

\section{Advanced Pulse Sequences}
\label{sec:e:advanced-pulse-sequences}
\subsection{Derivations and numerical checks}
\subsection{Interpretation, limiting cases and alternative approaches}

\section{Contrast and Relaxation}
\label{sec:e:contrast-and-relaxation}
\subsection{Derivations and numerical checks}
\subsection{Interpretation, limiting cases and alternative approaches}

\section{Diffusion MRI}
\label{sec:e:diffusion-mri}
\subsection{Derivations and numerical checks}
\subsection{Interpretation, limiting cases and alternative approaches}

\section{Quantitative MRI}
\label{sec:e:quantitative-mri}
\subsection{Derivations and numerical checks}
\subsection{Interpretation, limiting cases and alternative approaches}

\section{Magnetic Resonance Fingerprinting}
\label{sec:e:magnetic-resonance-fingerprinting}
\subsection{Derivations and numerical checks}
\subsection{Interpretation, limiting cases and alternative approaches}

\section{Extended Phase Graphs and Signal Pathways}
\label{sec:e:extended-phase-graphs-and-signal-pathways}
\subsection{Derivations and numerical checks}
\subsection{Interpretation, limiting cases and alternative approaches}

\section{Parallel Imaging}
\label{sec:e:parallel-imaging}
\subsection{Derivations and numerical checks}
\subsection{Interpretation, limiting cases and alternative approaches}

\section{Non-Cartesian and Accelerated Imaging}
\label{sec:e:non-cartesian-and-accelerated-imaging}
\subsection{Derivations and numerical checks}
\subsection{Interpretation, limiting cases and alternative approaches}

\section{Reconstruction and Inverse Problems}
\label{sec:e:reconstruction-and-inverse-problems}
\subsection{Derivations and numerical checks}
\subsection{Interpretation, limiting cases and alternative approaches}

\section{MRI Hardware}
\label{sec:e:mri-hardware}
\subsection{Derivations and numerical checks}
\subsection{Interpretation, limiting cases and alternative approaches}

\section{RF Pulses and Sequence Design}
\label{sec:e:rf-pulses-and-sequence-design}
\subsection{Derivations and numerical checks}
\subsection{Interpretation, limiting cases and alternative approaches}

\section{Flow, Perfusion and Angiography}
\label{sec:e:flow-perfusion-and-angiography}
\subsection{Derivations and numerical checks}
\subsection{Interpretation, limiting cases and alternative approaches}

\section{Magnetic Resonance Spectroscopy}
\label{sec:e:magnetic-resonance-spectroscopy}
\subsection{Derivations and numerical checks}
\subsection{Interpretation, limiting cases and alternative approaches}

\section{Functional MRI}
\label{sec:e:functional-mri}
\subsection{Derivations and numerical checks}
\subsection{Interpretation, limiting cases and alternative approaches}

\section{Artifacts and Corrections}
\label{sec:e:artifacts-and-corrections}
\subsection{Derivations and numerical checks}
\subsection{Interpretation, limiting cases and alternative approaches}

\section{Clinical Protocol Physics}
\label{sec:e:clinical-protocol-physics}
\subsection{Derivations and numerical checks}
\subsection{Interpretation, limiting cases and alternative approaches}

\section{Image Quality and Optimization}
\label{sec:e:image-quality-and-optimization}
\subsection{Derivations and numerical checks}
\subsection{Interpretation, limiting cases and alternative approaches}

\section{Modern MRI Methods}
\label{sec:e:modern-mri-methods}
\subsection{Derivations and numerical checks}
\subsection{Interpretation, limiting cases and alternative approaches}

\section{Safety and System Limits}
\label{sec:e:safety-and-system-limits}
\subsection{Derivations and numerical checks}
\subsection{Interpretation, limiting cases and alternative approaches}

\section{Integrated Worked Problems}
\label{sec:e:integrated-worked-problems}
\subsection{Derivations and numerical checks}
\subsection{Interpretation, limiting cases and alternative approaches}


% Index infrastructure is active; expand entries as chapters are drafted.
\cleardoublepage
\phantomsection
\addcontentsline{toc}{chapter}{Index}
\printindex
% Keep the compact reference sheet at the end of the book.
% Scope: End the book with a compact synthesis of established equations, conventions and parameter trade-offs; this pass creates headings only.
% Status: architecture only; no chapter prose or completed problems yet.
\chapter{Compact MRI Physics Reference Sheet}
\label{app:f}

\section{Spin and sequence essentials}
\label{sec:f:spin-and-sequence-essentials}
\subsection{Signs, units and rotations}
\subsection{Relaxation and sequence signal models}

\section{Encoding and reconstruction essentials}
\label{sec:f:encoding-and-reconstruction-essentials}
\subsection{Gradients, k-space and sampling}
\subsection{Coil encoding, noise and inverse problems}

\section{Quantification and design essentials}
\label{sec:f:quantification-and-design-essentials}
\subsection{Diffusion, flow and parameter estimation}
\subsection{Timing, bandwidth, SNR and system constraints}

\end{document}
